\chapter{The Spectral Flow \(N_{\mathrm{occ}}\rightsquigarrow H_{\mathrm{HP}}\): Operator and Domain Comparisons}
\label{v4-arith:w9-r4:hpar-spectral-flow}

The positive occupation spectrum and the zero-ordinate spectrum belong
to different operators. \Cref{v4-arith:w5-b:occupation-comparison}
constructs \(N_{\mathrm{occ}}v_m=(\log m)v_m\) and its exact embedding
in the product GNS space. The Hilbert--P\'olya target is a self-adjoint
operator whose shifted determinant has the actual zero divisor of
\(\xi_{\mathbb R}\). An intertwining problem must specify its carriers
and domains before a Mellin or residue formula can address it.

Two elementary obstructions are decisive. A bounded intertwiner from
the atomic occupation dynamics to the continuous Mellin dynamics is
zero. An injective intertwiner to a zero-divisor operator cannot
include the occupation vacuum in its domain. The continuous Mellin,
residue and rigged-space constructions below retain their own carriers;
comparison with the atomic arithmetic source is an additional problem.

\section{Setup: \(N_{\mathrm{occ}}\) and \(\Theta_{\xi}\) as Target Operators}
\label{v4-arith:w9-r4:setup}

The two endpoints of the spectral flow are the operators \(N_{\mathrm{occ}}\) and
\(H_{\mathrm{HP}}\), and the Hilbert spaces that carry them are specified below.

\subsection{Source: The Bost--Connes Log-Prime Operator}
\label{v4-arith:w9-r4:setup:source}

\begin{definition}[occupation Hilbert space and generator]
\label{v4-arith:w9-r4:def:bc-data}
Let \(\mathcal H_{\mathrm{BC}}=\ell^2(\mathbb N^\times)\), with
orthonormal basis \(v_m\), and let
\[
 N_{\mathrm{occ}}v_m=(\log m)v_m,\qquad
 \operatorname{Dom}N_{\mathrm{occ}}=
 \left\{(a_m):\sum_m(\log m)^2|a_m|^2<\infty\right\}.
\]
This is the self-adjoint occupation operator. At \(\beta=1\), the map
\(Wv_m=[\mu_m]\) is unitary onto the closed creation subspace
\(\mathcal H_1^+\) of the tensor-core GNS space and satisfies
\[
 WN_{\mathrm{occ}}=(-\log\Delta_1)W=2N_{\mathrm{BC}}W.
\]
Thus \(N_{\mathrm{occ}}\) is twice the restriction of the full
\(N_{\mathrm{BC}}=-\tfrac12\log\Delta_1\). The full space contains
division and matrix-unit sectors and has both signs of modular energy.
No global phase-algebra GNS identification is used here.
\end{definition}

\begin{remark}[weighted coordinates and prime factorization]
\label{v4-arith:w9-r4:rem:bc-structure}
For the weighted lattice \(\ell^2(\mathbb N^\times,dm/m)\), the
coordinate vector \(\delta_m\) has norm \(m^{-1/2}\). The isometries are
\(\delta_m\mapsto m^{-1/2}v_m\mapsto m^{-1/2}[\mu_m]\).
Prime factorization identifies the orthonormal basis \(v_m\) with the
finite-occupation tensor basis \(\bigotimes_p v_{v_p(m)}\), with
reference \(v_0\). The energy is \(\sum_p v_p(m)\log p\).
These Hilbert identifications concern the creation sector. It is not
invariant under modular conjugation or the full represented algebra.
At general \(\beta>0\), its modular eigenvalue is \(m^{-\beta}\),
so \(-\tfrac12\log\Delta_\beta\) restricts to
\((\beta/2)N_{\mathrm{occ}}\).
\end{remark}

\begin{proposition}[atomic-to-continuous intertwining obstruction]
\label{v4-arith:w9-r4:atomic-continuous-obstruction}
Let \(M_\tau\) be multiplication by \(\tau\) on \(L^2(\mathbb R,d\tau)\).
If a bounded linear map \(I:\mathcal H_{\mathrm{BC}}\to L^2(\mathbb R)\)
satisfies \(I e^{itN_{\mathrm{occ}}}=e^{itM_\tau}I\) for every real
\(t\), then \(I=0\). The same conclusion holds with any fixed Hilbert
space as multiplicity factor on the target.
\end{proposition}
\begin{proof}
Put \(f_m=Iv_m\). For every rational \(t\),
\((e^{it\tau}-e^{it\log m})f_m(\tau)=0\) almost everywhere.
Discard the union of the countably many exceptional null sets. If
\(\tau\ne\log m\), some rational \(t\) makes the parenthesis nonzero.
Thus \(f_m\) is supported at \(\{\log m\}\), a null set, and is zero.
Boundedness and density of the finite span of the \(v_m\)'s give
\(I=0\). The proof applies verbatim to Hilbert-valued functions.
\end{proof}

An embedding into distributions can send \(v_m\) to a Dirac mass at
\(\log m\), but a Dirac mass is not an \(L^2\) vector. Such a map
therefore requires a new topology and cannot supply a Hilbert-space
intertwiner without an additional comparison.

\subsection{Target: The Hypothetical Zeta-Zero Divisor Operator}
\label{v4-arith:w9-r4:setup:target}

\begin{definition}[target Hilbert space and zero-divisor operator]
\label{v4-arith:w9-r4:def:target-data}
\ClaimStatusDefinitional. Let \(\mathcal{Z}\) denote the multiset of
non-trivial zeros \(\rho=\beta_{\rho}+it_{\rho}\) of
\(\xi_{\mathbb{R}}\), with multiplicities \(m_{\rho}\), and write
\(\gamma_{\rho}:=t_{\rho}\) for the ordinate. The
\emph{formal target Hilbert space} is
\[
\mathcal{H}_{\xi}^{\mathrm{formal}}\;:=\;\bigoplus_{\rho\in\mathcal{Z}}\mathbb{C}^{m_{\rho}},
\]
with inner product making the one-dimensional summands
orthonormal. The \emph{formal ordinate multiplier} is the
diagonal symmetric operator
\[
H_{\mathrm{ord}}^{\mathrm{formal}}\;:\;(\xi_{\rho})_{\rho}\;\mapsto\;(t_{\rho}\xi_{\rho})_{\rho}
\]
on the finite-support core; its spectrum is
\(\{t_{\rho}\}\) with multiplicities. The \emph{critical-line
projection operator} is
\[
\Theta_{\xi}^{\mathrm{proj}}=\tfrac12\mathrm{id}+iH_{\mathrm{ord}}^{\mathrm{formal}}.
\]
Its divisor is the projection of the zero divisor to the critical
line; it is the actual zero divisor of \(\xi_{\mathbb R}\) exactly
under RH.
\end{definition}

\begin{remark}[why ``formal'']
\label{v4-arith:w9-r4:rem:why-formal}
The triple \((\mathcal{H}_{\xi}^{\mathrm{formal}},H_{\mathrm{ord}}^{\mathrm{formal}},\Theta_{\xi}^{\mathrm{proj}})\)
exists trivially as a diagonal construction. It carries no
intrinsic arithmetic structure: the enumeration of
\(\{t_{\rho}\}\) is read from the zero set itself, and the diagonal
operator realises its spectrum by fiat. If one replaces
\(\beta_{\rho}\) by \(1/2\), one has already projected the divisor to
the critical line. The content of the Hilbert--P\'olya problem is
not the existence of \(\mathcal{H}_{\xi}^{\mathrm{formal}}\); it is
the existence of a Hilbert space \emph{with independent structure}
whose spectral data force the original zero divisor to equal this
projected divisor.
\end{remark}

\begin{definition}[non-formal target]
\label{v4-arith:w9-r4:def:non-formal-target}
\ClaimStatusDefinitional. A \emph{non-formal target} is a pair
\((\mathcal{H}_{\xi},H_{\mathrm{HP}})\) where \(\mathcal{H}_{\xi}\) is a
Hilbert space carrying structure inherited from the Bost--Connes
framework (restricted-tensor adelic factorisation, Petersson inner
product on cuspidal cores, GNS modular data, or another
independently specified structure) and \(H_{\mathrm{HP}}\) is a
self-adjoint operator on \(\mathcal{H}_{\xi}\) whose spectral
resolution is canonical with respect to the inherited structure and
whose regularised determinant
\[
\det\nolimits_{\infty}\bigl(s\mathrm{id}-(\tfrac12\mathrm{id}+iH_{\mathrm{HP}})\bigr)
\]
has the non-trivial zero divisor of \(\xi_{\mathbb R}\).
\end{definition}

\begin{remark}[domain of ``canonical'']
\label{v4-arith:w9-r4:rem:canonical-domain}
``Canonical'' means: the spectral resolution of \(H_{\mathrm{HP}}\) is
determined by the inherited structure up to a unitary equivalence
whose parameters are themselves determined by that structure. A
diagonal construction with eigenvalues labelled by an arbitrary
enumeration is not canonical; a diagonal construction whose
eigenvectors are explicit functions of arithmetic data (adelic
characters, automorphic forms, spectral projectors of a globally
defined operator) is.
\end{remark}

\subsection{Spectral Flow as Intertwining Operator}
\label{v4-arith:w9-r4:setup:flow}

\begin{definition}[the spectral flow problem]
\label{v4-arith:w9-r4:def:spectral-flow}
\ClaimStatusDefinitional. The \emph{spectral flow problem} from
\(N_{\mathrm{occ}}\) to \(H_{\mathrm{HP}}\) is the specification of:
\begin{enumerate}
\item a Hilbert space \(\mathcal{H}_{\xi}\) with a non-formal structure
in the sense of \Cref{v4-arith:w9-r4:def:non-formal-target};
\item a densely defined linear operator
\(A:\mathcal{D}(A)\subset\mathcal{H}_{\mathrm{BC}}\to\mathcal{H}_{\xi}\)
with dense range and injective on its domain;
\item a self-adjoint operator \(H_{\mathrm{HP}}\) on \(\mathcal{H}_{\xi}\)
with domain \(\mathcal{D}(H_{\mathrm{HP}})\supset\mathrm{Ran}(A)\);
\item the intertwining identity
\begin{equation}
\label{v4-arith:w9-r4:eq:intertwine}
A\,N_{\mathrm{occ}}\,A^{-1}\;=\;H_{\mathrm{HP}}\qquad\text{on }\mathrm{Ran}(A)\cap\mathcal{D}(H_{\mathrm{HP}}),
\end{equation}
interpreted in the closed-operator sense after specifying domain
compatibility;
\item the spectral identity
\begin{equation}
\label{v4-arith:w9-r4:eq:spec-id}
\mathrm{Spec}(H_{\mathrm{HP}})\;=\;\{\gamma_{\rho}\mid\rho\in\mathcal{Z}\}
\end{equation}
with multiplicities matching.
\end{enumerate}
A \emph{solution} is a tuple \((\mathcal{H}_{\xi},A,H_{\mathrm{HP}})\)
satisfying all five items. The \emph{Hilbert--P\'olya twist} is a
solution.
\end{definition}

\begin{proposition}[vacuum obstruction to literal spectral flow]
\label{v4-arith:w9-r4:vacuum-flow-obstruction}
The spectral-flow conditions cannot hold if \(v_1\in\mathcal D(A)\)
and the intertwining identity applies there, with a target whose
shifted determinant has the actual divisor of \(\xi_{\mathbb R}\).
\end{proposition}
\begin{proof}
Since \(N_{\mathrm{occ}}v_1=0\), intertwining gives
\(H_{\mathrm{HP}}Av_1=0\). Injectivity gives \(Av_1\ne0\).
The shifted operator then has eigenvalue \(1/2\), so its determinant
divisor would contain \(1/2\). But \(\xi_{\mathbb R}(1/2)\ne0\), as
proved in \Cref{v4-arith:w5-b:vacuum-obstruction}.
\end{proof}

Any remaining literal spectral-flow variant must specify a domain
excluding this test or an altered comparison. A construction with a
separate occupation grading and zero operator, as in
\Cref{v4-arith:w5-b:def:twist}, does not assert literal intertwining.

\begin{remark}[unitarity of \(A\)]
\label{v4-arith:w9-r4:rem:unitarity-of-A}
\Cref{v4-arith:w9-r4:def:spectral-flow} does \emph{not} require \(A\)
to be bounded or unitary. The closed-operator intertwining
\eqref{v4-arith:w9-r4:eq:intertwine} admits unbounded \(A\); in the
regular case one asks for a polar decomposition \(A=U|A|\) with \(U\)
partial isometry and \(|A|\) self-adjoint non-negative, on which the
intertwining reads
\(U|A|N_{\mathrm{occ}}|A|^{-1}U^{*}=H_{\mathrm{HP}}\) on its natural
domain.  It becomes a unitary intertwining of one-parameter groups only
when the polar partial isometry is unitary between the two Hilbert
closures, the domains are invariant under the two Stone groups, and the
closed-operator identity holds on a common core.  Then
\(U\,\exp(it N_{\mathrm{occ}})\,U^{*}=\exp(itH_{\mathrm{HP}})\). We use
the closed-operator form throughout; when the literature writes
``\(A\) is a unitary twist'' it is the polar \(U\) in this restricted
sense.
\end{remark}

\section{Prime-Zero Mangoldt-Hadamard Duality Recap}
\label{v4-arith:w9-r4:mangoldt-hadamard}

The formal bridge between the two spectra is the Riemann--von Mangoldt
explicit formula and its Hadamard factorization. The prime measure acts on compact smooth
tests; its Fourier use requires the explicit exponential weights
and entire test carrier specified below. These identities are
distributional pairings and do not construct a Hilbert-space operator.

\subsection{Hadamard Factorisation of \(\xi_{\mathbb{R}}\)}
\label{v4-arith:w9-r4:mh:hadamard}

\begin{proposition}[Hadamard factorisation]
\label{v4-arith:w9-r4:prop:hadamard}
\ClaimStatusProvedElsewhere \emph{(Hadamard 1893;
\Cref{v4-arith:w5-b:eq:hadamard})}. The completed Riemann zeta
\(\xi_{\mathbb{R}}(s)=\tfrac12 s(s-1)\pi^{-s/2}\Gamma(s/2)\zeta(s)\) is
entire of order \(1\), of genus \(1\), and admits the factorisation
\begin{equation}
\label{v4-arith:w9-r4:eq:hadamard-factor}
\xi_{\mathbb{R}}(s)\;=\;\xi_{\mathbb{R}}(0)\,e^{Bs}\,\prod_{\rho}\Bigl(1-\frac{s}{\rho}\Bigr)e^{s/\rho},
\end{equation}
with \(\rho\) ranging over non-trivial zeros with multiplicity and
\[
 B=\frac{\xi_{\mathbb R}'(0)}{\xi_{\mathbb R}(0)}
   =-1-\frac\gamma2+\frac12\log(4\pi).
\]
The product uses genus-one factors. Its convergence and this exact
normalization are proved in Proposition~\ref{v4-arith:pm:constant}.
\end{proposition}

\subsection{Riemann-von Mangoldt Explicit Formula}
\label{v4-arith:w9-r4:mh:rvm}

\begin{proposition}[Riemann-von Mangoldt explicit formula]
\label{v4-arith:w9-r4:prop:rvm}
\ClaimStatusProvedElsewhere \emph{(von Mangoldt 1895; Edwards 1974 \S3)}.
For \(x>1\) non-integer,
\begin{equation}
\label{v4-arith:w9-r4:eq:rvm}
\sum_{n\leq x}\Lambda(n)\;=\;x\;-\;\sum_{\rho}\frac{x^{\rho}}{\rho}\;-\;\frac{\zeta'(0)}{\zeta(0)}\;-\;\tfrac12\log(1-x^{-2}),
\end{equation}
where \(\Lambda(n)=\log p\) if \(n=p^{k}\) and \(0\) otherwise.
The zero sum is the symmetric height limit
$\lim_{T\to\infty}\sum_{|\Im\rho|\leq T}x^\rho/\rho$.
At a prime power the left side is replaced by the average of its
left and right limits. The constant is
$\zeta'(0)/\zeta(0)=\log(2\pi)$, distinct from the Hadamard constant
$B$ of the preceding product. The smooth-test formulation, including
all convergence and Fourier normalizations, is proved below.
\end{proposition}

\subsection{Normalization and the spaces of prime distributions}
\label{v4-arith:pm:chapter}

The exponential coordinate $m=e^t$ changes the growth of a counting
measure. It also fixes the weight appearing in a Fourier formula.
We use additive Dirac masses on $\mathbb R$, the test space
$\mathscr A=C_c^\infty(\mathbb R)$ with its usual inductive-limit
topology, and the Schwartz space $\mathscr S(\mathbb R)$.
Distributions are complex-linear in their test functions.
Write $\gamma$ for Euler's constant and
\[
 \xi(s)=\xi_{\mathbb R}(s)
   =\tfrac12s(s-1)\pi^{-s/2}\Gamma(s/2)\zeta(s).
\]

\begin{lemma}[growth and zero counting from the theta integral]
\label{v4-arith:pm:zero-count}
For $R\geq2$ one has
$\log\max_{|s|\leq R}|\xi(s)|\leq CR\log(R+2)$.
Consequently the number of its zeros in $|s|\leq R$, with
multiplicity, is $O(R\log(R+2))$, and
$\sum_\rho|\rho|^{-2}<\infty$.
There are infinitely many zeros. All lie in $0\leq\Re s\leq1$.
\end{lemma}
\begin{proof}
In the theta integral \eqref{v4-arith:mo:xi-integral}, the theta
tail on $x\geq1$ is bounded by $C e^{-\pi x}$. For $|s|\leq R$
the powers of $x$ in that integral are bounded by $2x^{(R+1)/2}$.
The estimate
$\sup_{x>0}x^A e^{-\pi x/2}=(2A/(\pi e))^A$
bounds the integral by $\exp(CR\log(R+2))$; the quadratic
prefactor does not change this bound. Jensen's formula at radius
$2R$, using $\xi(0)=1/2$, bounds each zero in radius $R$ by a
contribution at least $\log2$. This proves the zero count.
Dyadic summation then proves the inverse-square assertion.

The growth bound gives entire order at most one. If there were
only finitely many zeros, Hadamard's factorization theorem would
give $\xi(s)=e^{a+bs}P(s)$ with a polynomial $P$.
For real $s\geq2$, however, $\zeta(s)\geq1$ and integration of
the gamma integral just over $[s/2,s/2+1]$ gives
$\log\xi(s)\geq(s/2)\log s-Cs$. This contradicts that finite
factorization. Finally, the Euler product excludes zeros for
$\Re s>1$, and the functional equation excludes $\Re s<0$.
\end{proof}

\begin{proposition}[the constant in the genus-one product]
\label{v4-arith:pm:constant}
For this normalization, $\xi(0)=1/2$, and the constant multiplying
$s$ in the canonical genus-one product is
\[
 B=\frac{\xi'(0)}{\xi(0)}
   =-1-\frac\gamma2+\frac12\log(4\pi).
\]
Thus, with zeros counted with multiplicity,
\[
 \xi(s)=\tfrac12e^{Bs}\prod_\rho(1-s/\rho)e^{s/\rho},\qquad
 \frac{\xi'(s)}{\xi(s)}
   =B+\sum_\rho\left(\frac1{s-\rho}+\frac1\rho\right).
\]
The factors and the paired summands converge locally uniformly away
from their indicated zeros or poles. The individual sum
$\sum_\rho1/\rho$ is not separated from its partner without a
specified summation convention.
\end{proposition}
\begin{proof}
The identity $\Gamma(1+s/2)=(s/2)\Gamma(s/2)$ removes the apparent
singularity at zero:
\[
 \xi(s)=(s-1)\Gamma(1+s/2)\pi^{-s/2}\zeta(s).
\]
The values $\zeta(0)=-1/2$, $\zeta'(0)=-\tfrac12\log(2\pi)$,
and $\Gamma'(1)=-\gamma$ give
\[
 \frac{\xi'(0)}{\xi(0)}
 =-1-\frac\gamma2-\frac12\log\pi+\log(2\pi).
\]
For example, the two zeta values follow directly by expanding
$\zeta(s)=2^s\pi^{s-1}\sin(\pi s/2)\Gamma(1-s)\zeta(1-s)$:
the constant terms cancel in
$\Gamma(1-s)\zeta(1-s)=-s^{-1}+O(s)$.
The remaining factor is
$\tfrac s2(1+s\log(2\pi)+O(s^2))$.
The exact special values are also recorded in
\href{https://dlmf.nist.gov/25.6.E1}{DLMF 25.6.1},
\href{https://dlmf.nist.gov/25.6.E11}{25.6.11}, and
\href{https://dlmf.nist.gov/5.4.E11}{5.4.11}.

Lemma~\ref{v4-arith:pm:zero-count} supplies the growth and zero bound
for the order-one factorization. On a compact set with
$|s/\rho|\leq1/2$, the logarithm of a factor and its logarithmic
derivative are bounded by a constant times $|\rho|^{-2}$.
Hence the product and paired derivative converge normally there.
Every canonical factor has derivative zero at $s=0$, so its
exponential constant is precisely the logarithmic derivative just
computed. This also specifies the normalization of the product.
\end{proof}

\begin{lemma}[an elementary lower bound for logarithmic prime mass]
\label{v4-arith:pm:growth}
Put $\psi(x)=\sum_{m\leq x}\Lambda(m)$. For every integer $n\geq1$,
\[
 2n\log2-\log(2n+1)\leq\psi(2n),
 \qquad \psi(x)\leq x\log x\quad(x\geq2).
\]
Consequently there are positive constants $c,x_0$ with
$\psi(x)\geq cx$ for $x\geq x_0$.
\end{lemma}
\begin{proof}
The largest of the $2n+1$ binomial coefficients in $(1+1)^{2n}$
is $\binom{2n}{n}$, so it is at least $4^n/(2n+1)$.
For every prime $p$, its valuation is
\[
 v_p\binom{2n}{n}
 =\sum_{k\geq1}\left(\left\lfloor\frac{2n}{p^k}\right\rfloor
                 -2\left\lfloor\frac n{p^k}\right\rfloor\right).
\]
Each summand is zero or one and vanishes when $p^k>2n$.
Taking logarithms therefore gives
$\log\binom{2n}{n}\leq\psi(2n)$.
The upper bound uses $0\leq\Lambda(m)\leq\log m$.
For arbitrary sufficiently large $x$, take $n=\lfloor x/2\rfloor$
and use monotonicity. No prime-number asymptotic is needed.
\end{proof}

\begin{theorem}[the exact exponential weight for temperedness]
\label{v4-arith:pm:weights}
For real $\sigma$, define the positive Radon measure
\[
 \mu_\sigma=\sum_{m\geq2}\Lambda(m)m^{-\sigma}\delta_{\log m}.
\]
Every $\mu_\sigma$ is a distribution on $\mathscr A$.
It is a tempered distribution exactly when $\sigma\geq1$.
It is finite exactly when $\sigma>1$; $\mu_1$ has infinite total mass.
In particular, $\mu_0=\mu_{\mathrm{prime}}$ is unchanged as a
locally finite measure, and $\mu_{1/2}$ is not tempered.
\end{theorem}
\begin{proof}
Only finitely many atoms meet any compact interval. Their total
variation bounds evaluation by the supremum norm on that interval,
which proves continuity on $\mathscr A$.
For $\sigma\geq1$,
\[
 |\mu_\sigma(\varphi)|
 \leq\sup_t(1+|t|)^3|\varphi(t)|
       \sum_{m\geq2}\frac{\log m}{m(1+\log m)^3}.
\]
The last series converges by the integral test. If $\sigma>1$,
$\sum(\log m)m^{-\sigma}<\infty$ also gives finiteness.
Partial summation gives
\[
 \sum_{m\leq X}\frac{\Lambda(m)}m
 =\frac{\psi(X)}X+\int_1^X\frac{\psi(x)}{x^2}\,dx,
\]
which diverges by Lemma~\ref{v4-arith:pm:growth}.

If a positive measure supported in $[0,\infty)$ is tempered, choose
a nonnegative smooth compact cutoff equal to one on $[0,1]$ and
rescale it by $t\mapsto t/T$. The finite Schwartz seminorm bound
for a tempered distribution then implies
$\mu([0,T])\leq C(1+T)^M$ for some $C,M$ and all $T\geq1$.
For $0\leq\sigma<1$, however,
\[
 \mu_\sigma([0,T])\geq e^{-\sigma T}\psi(e^T)
                  \geq c e^{(1-\sigma)T}
\]
for large $T$. For $\sigma<0$, use $\mu_\sigma\geq\mu_0$.
Both contradict polynomial mass growth.
\end{proof}

Multiplication by $e^{-\sigma t}$ is explicit here; it is not a
coordinate relabelling. Equivalently, for $\sigma\geq1$ the
unweighted $\mu_0$ acts continuously on
$\mathscr S_\sigma=e^{-\sigma t}\mathscr S(\mathbb R)$ with the
transported Schwartz topology. It acts on $C_c^\infty$ without any
choice of weight. The change $x=e^t$ pulls a compact smooth test
back by $f(t)=g(e^t)$ and gives the exact pairing
$\mu_\sigma(f)=\sum\Lambda(m)m^{-\sigma}g(m)$.
These are maps between specified test spaces, not Hilbert-space
intertwiners.

\begin{proposition}[Laplace and Fourier transforms on their domains]
\label{v4-arith:pm:laplace}
For $\Re s>1$,
\[
 \int e^{-st}\,d\mu_0(t)
 =\sum_{m\geq2}\frac{\Lambda(m)}{m^s}
 =-\frac{\zeta'(s)}{\zeta(s)}.
\]
With Fourier convention $\mathcal F_-f(u)=\int f(t)e^{-iut}dt$,
the finite measure $\mu_\sigma$, $\sigma>1$, has Fourier transform
$-\zeta'(\sigma+iu)/\zeta(\sigma+iu)$.
As $\sigma\downarrow1$ these transforms converge in the weak
tempered-distribution sense to $\mathcal F_-\mu_1$.
The unweighted and critically weighted measures do not have Fourier
transforms in $\mathscr S'$ by this construction.
\end{proposition}
\begin{proof}
The absolutely convergent Euler product for $\Re s>1$ may be
logarithmically differentiated on compact subsets there, giving
the displayed Dirichlet series. Absolute convergence also permits
integration of the finite weighted measure against the Fourier
kernel. For every Schwartz test, dominated convergence in the
seminorm estimate of Theorem~\ref{v4-arith:pm:weights} gives
$\mu_{1+\varepsilon}\to\mu_1$ weakly in $\mathscr S'$.
Fourier continuity proves the boundary assertion. Meromorphic
continuation of the right side does not change the abscissa of
absolute convergence of the original positive-measure integral.
\end{proof}

\subsection{The explicit formula on compact Fourier tests}
\label{v4-arith:pm:explicit-section}

For $f\in\mathscr A$, use the entire function
$h(z)=\int f(t)e^{izt}dt$, so that
$f(t)=(2\pi)^{-1}\int h(u)e^{-iut}du$.
This is the Fourier test carrier of
Definition~\ref{v4-arith:mc:test-algebra}.
For $\rho=\beta+i\gamma_\rho$ put
$z_\rho=-i(\rho-1/2)=\gamma_\rho-i(\beta-1/2)$.
Lemma~\ref{v4-arith:pm:zero-count} gives
$\#\{\rho:|\gamma_\rho|\leq T\}=O(T\log(T+2))$, with multiplicity.
It follows that $\mu_{\mathrm{ord}}=\sum_\rho\delta_{\gamma_\rho}$
is tempered. Since $0\leq\beta\leq1$, integration by parts gives rapid
decay of $h(z_\rho)$ uniformly in this horizontal strip.
Thus $\mathfrak Z_\xi(h)=\sum_\rho h(z_\rho)$ converges absolutely
and is continuous on $\mathscr A$ through $f\mapsto h$.
The estimates therefore follow from the theta integral without
requiring an asymptotic prime-counting formula.

\begin{theorem}[the full polar and archimedean terms]
\label{v4-arith:pm:weil}
For every $f\in C_c^\infty(\mathbb R)$ and its transform $h$ above,
\begin{align*}
 \mathfrak Z_\xi(h)
 &=h(i/2)+h(-i/2)
   +\frac1{2\pi}\int_{\mathbb R}
       h(u)\bigl(\Re\psi_\Gamma(1/4+iu/2)-\log\pi\bigr)\,du\\
 &\quad-\sum_{m\geq2}\frac{\Lambda(m)}{\sqrt m}
                  \bigl(f(\log m)+f(-\log m)\bigr),
\end{align*}
where $\psi_\Gamma=\Gamma'/\Gamma$ is the digamma function.
The prime sum is finite. In particular, it is evaluation of
$\mu_{1/2}$ on compact tests and does not assert temperedness of
that measure. The values $h(\pm i/2)$ are the polar terms, separate
from the real-line digamma density.
\end{theorem}
\begin{proof}
Set $H(s)=h(-i(s-1/2))=\int f(t)e^{(s-1/2)t}dt$ and choose $a>1$.
The residue theorem and $\xi(s)=\xi(1-s)$ give
\[
 \sum_\rho H(\rho)
 =\frac1{2\pi i}\int_{\Re s=a}
             (H(s)+H(1-s))\frac{\xi'(s)}{\xi(s)}\,ds.
\]
Here is a convergence justification. The zero count allows heights
$T_j\in[j,j+1]$ at distance at least $j^{-2}$ from every zero
ordinate: the excluded intervals have total length $O(\log j/j)$.
The paired logarithmic derivative in
Proposition~\ref{v4-arith:pm:constant} is polynomially bounded on
the horizontal segments at these heights. For $|\rho|\leq2T_j$
bound each denominator by the stated distance and sum
$O(T_j\log T_j)$ terms; for $|\rho|>2T_j$ use
$|s/\{\rho(s-\rho)\}|\leq C T_j/|\rho|^2$.
Repeated integration by parts makes $H$ decrease faster than any
power, uniformly on the fixed strip. The horizontal integrals
therefore vanish, and the zero sum converges absolutely.

On the right line decompose
\[
 \frac{\xi'}\xi(s)
 =\frac1s+\frac1{s-1}-\frac12\log\pi
       +\frac12\psi_\Gamma(s/2)+\frac{\zeta'}\zeta(s).
\]
Fourier inversion and the absolutely convergent Dirichlet series
give the negative prime sum from the last term. For the remaining
terms move the line to $\Re s=1/2$. The crossed pole at $s=1$
contributes $H(1)+H(0)$. On the new line the rational sum
$1/s+1/(s-1)$ is odd in the imaginary coordinate, while
$H(s)+H(1-s)$ is even; its integral is zero. Changing $u$ to $-u$
in half of the gamma integral gives precisely its real part in the
displayed formula. Digamma growth is $O(\log(|u|+2))$ on this
line: in
$\psi_\Gamma(z)=-\gamma+\sum_{n\geq0}((n+1)^{-1}-(n+z)^{-1})$,
split the sum at $n=|u|+1$ for $z=1/4+iu/2$. The initial sum
is $O(\log(|u|+2))$ and the tail is bounded by
$C(|u|+1)\sum_{n>|u|+1}n^{-2}=O(1)$.
Thus all resulting integrals converge absolutely.
\end{proof}

When $f$ is even, $h$ is even and the prime term is
$-2\mu_{1/2}(f)$. The functional
$h\mapsto h(i/2)+h(-i/2)$ cannot be silently included in a density
on the real line. Likewise, $\mathfrak Z_\xi(h)$ uses the full
complex numbers $z_\rho$, whereas
$\mu_{\mathrm{ord}}(h)=\sum h(\gamma_\rho)$ uses their real parts.
Under RH these are equal. No projection is used in the theorem.

\subsection{Centering and the critical exponent}
\label{v4-arith:pm:centered-section}

For real $\sigma$ define, initially on compact smooth tests,
\[
 \nu_\sigma=\mu_\sigma
       -\mathbf1_{[0,\infty)}(t)e^{(1-\sigma)t}\,dt.
\]
This subtracts an explicit density and is a different distribution.
For $\sigma\geq1$ it is tempered by
Theorem~\ref{v4-arith:pm:weights}. In particular its transform at
$\sigma=1$ is the weak boundary value
\[
 \mathcal F_-\nu_1
 =\lim_{\varepsilon\downarrow0}
    \left(-\frac{\zeta'(1+\varepsilon+iu)}
                 {\zeta(1+\varepsilon+iu)}
                    -\frac1{\varepsilon+iu}\right).
\]
The removed half-line density has Fourier transform
$\pi\delta_0-i\,\mathrm{pv}(1/u)$, by taking the real and imaginary
parts of $(\varepsilon+iu)^{-1}$. Thus centering specifies both the
Dirac and principal-value boundary terms.

\begin{theorem}[critical centering as a distributional criterion]
\label{v4-arith:pm:centered}
The distribution $\nu_{1/2}$ is tempered if and only if every
nontrivial zero of $\zeta$ has real part $1/2$.
This is a criterion for a specified centered distribution; it
constructs no arithmetic Hilbert pairing or occupation intertwiner.
\end{theorem}
\begin{proof}
First restrict a test $f$ to $C_c^\infty(0,\infty)$.
The digamma integral representation
\[
 \psi_\Gamma(z)=-\gamma+
  \int_0^\infty\frac{e^{-v}-e^{-zv}}{1-e^{-v}}\,dv
 \qquad(\Re z>0)
\]
is \href{https://dlmf.nist.gov/5.9.E16}{DLMF 5.9.16}.
This identity
and Fourier inversion, with $v=2t$, turn the archimedean integral
of Theorem~\ref{v4-arith:pm:weil} into
$-\int_0^\infty f(t)e^{-t/2}/(1-e^{-2t})\,dt$.
The constant terms vanish since $f(0)=0$; near $v=0$ the
subtracted representation is integrable, and compact support away
from zero justifies the Fourier inversion. The term $h(i/2)$
cancels the first exponential in this denominator. Hence, on
$(0,\infty)$,
\[
 \nu_{1/2}(f)
 =-\sum_\rho\int f(t)e^{(\rho-1/2)t}\,dt
   -\int_0^\infty\frac{e^{-t/2}}{e^{2t}-1}f(t)\,dt.
\]
Under RH the zero sum extends to a tempered distribution on the
whole real line by the exact map
$f\mapsto\mu_{\mathrm{ord}}(\mathcal F_+f)$, where
$\mathcal F_+f(u)=\int f(t)e^{iut}dt$.
Multiply this distribution by a smooth cutoff that vanishes near
zero and is one for $t\geq1$. The last displayed density then
decays exponentially. The difference from $\nu_{1/2}$ has compact
support and is a distribution, so it too is tempered. This proves
the forward construction under RH.

Conversely suppose $\nu_{1/2}\in\mathscr S'$. Its support is
contained in $[0,\infty)$. Choose a smooth function $\eta$ that
is one there and zero on $(-\infty,-1]$. For $\Re s>0$ define
$L(s)=\nu_{1/2}(\eta(t)e^{-st})$.
The tests and all their $s$-derivatives lie in $\mathscr S$, locally
uniformly for $\Re s>0$; hence $L$ is holomorphic there.
For $\Re s>1/2$ the convergent integral gives
\[
 L(s)=-\frac{\zeta'(s+1/2)}{\zeta(s+1/2)}
                         -\frac1{s-1/2}.
\]
Meromorphic uniqueness extends this equality to $\Re s>0$.
The pole at $s=1/2$ cancels, but any zero with $\Re\rho>1/2$
would leave a pole of nonzero integer residue at $s=\rho-1/2$.
Holomorphy excludes these zeros. The functional equation excludes
zeros with $\Re\rho<1/2$ as well.
\end{proof}

The unweighted centered distribution $\nu_0$ is not tempered.
Indeed the same Laplace argument would make
$-\zeta'(s)/\zeta(s)-1/(s-1)$ holomorphic for $\Re s>0$.
The infinitely many xi zeros of Lemma~\ref{v4-arith:pm:zero-count},
together with reflection in $\Re s=1/2$, provide a zero with
positive real part, and hence an uncancelled pole.

\begin{corollary}[the ordinate shadow on the full test class]
\label{v4-arith:pm:shadow}
Equality $\mathfrak Z_\xi(h)=\mu_{\mathrm{ord}}(h)$ for every
$h(z)=\int f(t)e^{izt}dt$, $f\in C_c^\infty(\mathbb R)$,
is equivalent to RH.
\end{corollary}
\begin{proof}
Under RH the equality is termwise. Conversely the equality makes
the zero functional on compact Fourier tests the restriction of a
tempered Fourier transform. The cutoff argument in
Theorem~\ref{v4-arith:pm:centered} then proves
$\nu_{1/2}\in\mathscr S'$. That theorem implies RH.
\end{proof}

The uncentered critical measure remains non-tempered even under RH.
Centering subtracts its exponentially growing main density; a
statement about this difference does not change the carrier of the
original prime measure. These distinctions are required before
any Fourier, Mellin, or spectral comparison can be made.


\subsection{Mellin-Dual Measures}
\label{v4-arith:w9-r4:mh:mellin-dual}

\begin{definition}[prime measure and full zero functional]
\label{v4-arith:w9-r4:def:counting-measures}
\ClaimStatusDefinitional. The unweighted prime measure is
\begin{equation}
\label{v4-arith:w9-r4:eq:mu-prime}
 \mu_{\mathrm{prime}}=\mu_0
 =\sum_{m\geq2}\Lambda(m)\delta_{\log m}
 =\sum_p\sum_{k\geq1}(\log p)\delta_{k\log p}.
\end{equation}
It is a positive Radon measure on $[0,\infty)$ and a distribution
on $C_c^\infty(\mathbb R)$, but it is not tempered.
The weighted measure is exactly $\mu_\sigma=e^{-\sigma t}\mu_0$;
Theorem~\ref{v4-arith:pm:weights} gives its temperedness threshold.
The ordinate shadow is the tempered measure
\begin{equation}
\label{v4-arith:w9-r4:eq:mu-zero}
 \mu_{\mathrm{ord}}=\sum_\rho\delta_{\gamma_\rho}.
\end{equation}
On the Fourier images $h(z)=\int f(t)e^{izt}dt$ of compact smooth
tests, the actual zero functional is the absolutely convergent sum
\[
 \mathfrak Z_\xi(h)=\sum_\rho h(-i(\rho-1/2)).
\]
Its equality with $\mu_{\mathrm{ord}}(h)$ on this entire test
class is equivalent to RH by Corollary~\ref{v4-arith:pm:shadow}.
\end{definition}

\begin{proposition}[the weighted prime term in the explicit formula]
\label{v4-arith:w9-r4:prop:mellin-dual}
\ClaimStatusProvedHere. Let $f\in C_c^\infty(\mathbb R)$ be even
and $h(z)=\int f(t)e^{izt}dt$. Then
\begin{equation}
\label{v4-arith:w9-r4:eq:mellin-dual}
 \mathfrak Z_\xi(h)=h(i/2)+h(-i/2)+\mathcal A_\infty(h)
                         -2\mu_{1/2}(f),
\end{equation}
where
\[
 \mathcal A_\infty(h)=\frac1{2\pi}\int_{\mathbb R}
 h(u)\bigl(\Re\psi_\Gamma(1/4+iu/2)-\log\pi\bigr)\,du.
\]
Here $f(t)=(2\pi)^{-1}\int h(u)e^{-iut}du$ is the inverse Fourier
transform, with its normalization fixed. The two polar evaluations
are separate from the real-line archimedean density.
\end{proposition}
\begin{proof}
Apply Theorem~\ref{v4-arith:pm:weil} and use evenness of $f$.
Only finitely many prime powers meet its support. The prime term
therefore evaluates a locally finite measure on a compact test;
it uses no tempered Fourier transform of $\mu_{1/2}$.
\end{proof}

\begin{remark}[distributional pairings and operator comparisons]
\label{v4-arith:w9-r4:rem:duality-distributional}
The explicit formula compares the full zero functional on an
entire Fourier test class with a compact-test pairing against
$e^{-t/2}\mu_{\mathrm{prime}}$, two polar evaluations, and the
digamma density. It does not identify the unweighted measure with
a tempered distribution. Weighting by $e^{-t}$ does give a
tempered measure; centering at exponent $1/2$ gives a different
distribution whose temperedness is the criterion of
Theorem~\ref{v4-arith:pm:centered}. Neither operation is implicit
in the formula. A map between Hilbert spaces and their operator
domains still requires the separate constructions below.
\end{remark}

\section{Formal Spectral Flow \(A\) as Mellin-Inversion Symbol}
\label{v4-arith:w9-r4:formal-flow}

The candidate spectral-flow operator exists at the symbolic level as
the composite of Mellin transform and Hadamard residue extraction;
the two obstructions that prevent its promotion to a Hilbert-space
operator are named below.

\subsection{The Mellin-Hadamard Inversion Symbol}
\label{v4-arith:w9-r4:formal-flow:symbol}

\begin{definition}[formal Mellin-Hadamard inversion]
\label{v4-arith:w9-r4:def:formal-MH}
\ClaimStatusDefinitional. The \emph{formal Mellin-Hadamard inversion}
is the composite symbol
\begin{equation}
\label{v4-arith:w9-r4:eq:formal-MH}
A^{\mathrm{sym}}\;:=\;\mathcal{I}_{\mathrm{Had}}\circ\mathcal{M}
\end{equation}
where \(\mathcal{M}\) is the Mellin transform
\((\mathcal{M}f)(s)=\int_{0}^{\infty}f(x)x^{s-1}dx\) on a suitable
function class and \(\mathcal{I}_{\mathrm{Had}}\) is the
Hadamard-inversion formal symbol
\[
(\mathcal{I}_{\mathrm{Had}}F)(\gamma_{\rho})\;:=\;\operatorname*{Res}_{s=\rho}\frac{F(s)}{\xi_{\mathbb{R}}(s)}
\]
reading the residue of \(F/\xi_{\mathbb{R}}\) at each non-trivial zero
and writing it into the \(\rho\)-component of a formal sequence
indexed by \(\mathcal{Z}\).
\end{definition}

\begin{remark}[domains of the residue map]
\label{v4-arith:w9-r4:rem:Asym-not-operator}
On $C_c^\infty(\mathbb R_{>0})$ the displayed formula is a linear
map to all zero-indexed scalar sequences. On finite occupation
sequences it becomes a linear map after the distributional embedding
of Proposition~\ref{v4-arith:mo:occupation} has been specified.
A Hilbert-space operator additionally requires a target norm,
convergence of the full residue sequence, and a source domain.
Closedness must be checked in those two norms.
Theorem~\ref{v4-arith:mo:finite-graph} and Theorem~\ref{v4-arith:mo:all-jets} construct
evaluation completions and compute their graphs explicitly.
\end{remark}

\subsection{Formal Mellin--Hadamard Obstruction}
\label{v4-arith:w9-r4:formal-flow:formal-intertwine}

\begin{lemma}[formal Mellin--Hadamard obstruction]
\label{v4-arith:w9-r4:lem:formal-intertwine}
\ClaimStatusProvedHere. On \(C_c^\infty(\mathbb R_{>0})\), define
\((M_{\log}f)(x)=-(\log x)f(x)\). This continuous multiplication
operator is distinct from the atomic occupation operator. Suppose \(f\in C_{c}^{\infty}(\mathbb{R}_{>0})\)
has Mellin transform \(F(s)=(\mathcal{M}f)(s)\) entire, and suppose
\(\rho\in\mathcal{Z}\) is a simple zero of \(\xi_{\mathbb R}\). Then
\begin{equation}
\label{v4-arith:w9-r4:eq:formal-intertwine}
A^{\mathrm{sym}}\bigl(M_{\log}f\bigr)(\gamma_{\rho})
\;=\;
-\,\frac{F'(\rho)}{\xi_{\mathbb R}'(\rho)},
\qquad
A^{\mathrm{sym}}(f)(\gamma_{\rho})
\;=\;
\frac{F(\rho)}{\xi_{\mathbb R}'(\rho)}.
\end{equation}
On formal sequences indexed by the zeros, define
$(M_\gamma a)_\rho=\gamma_\rho a_\rho$.
Consequently the bare residue intertwining relation
\[
A^{\mathrm{sym}}M_{\log}
\stackrel{?}{=}
M_{\gamma}A^{\mathrm{sym}}
\]
holds at \(\rho\) only under the additional interpolation condition
\[
F'(\rho)=-\gamma_{\rho}F(\rho).
\]
This condition is not a consequence of RH.
\end{lemma}

\begin{proof}
Differentiate the Mellin integral of the compactly supported function.
The Mellin transform gives
\((\mathcal{M}(M_{\log}f))(s)=-F'(s)\), where \(F'\) is the
derivative in the spectral variable. Since \(\rho\) is a simple zero,
\[
\operatorname*{Res}_{s=\rho}\frac{F(s)}{\xi_{\mathbb R}(s)}
=\frac{F(\rho)}{\xi_{\mathbb R}'(\rho)},
\qquad
\operatorname*{Res}_{s=\rho}\frac{-F'(s)}{\xi_{\mathbb R}(s)}
=-\frac{F'(\rho)}{\xi_{\mathbb R}'(\rho)}.
\]
Multiplication by the zero ordinate would instead give
\(\gamma_{\rho}F(\rho)/\xi_{\mathbb R}'(\rho)\). Equality is exactly
the displayed first-order interpolation condition.
\end{proof}

\begin{remark}[RH is not enough for formal intertwining]
\label{v4-arith:w9-r4:rem:formal-assumes-RH}
\Cref{v4-arith:w9-r4:lem:formal-intertwine} shows that multiplication by $-\log x$
differentiates the numerator before Hadamard residue
extraction. RH would place \(\rho\) on the critical line, but it would
not imply \(F'(\rho)=-\gamma_{\rho}F(\rho)\) for the Mellin transforms
in the source class. Thus the bare symbol
\(A^{\mathrm{sym}}\) is not a spectral-flow operator. A genuine
operator construction must supply a new domain theorem or boundary
condition that converts derivative residue data into zero-ordinate
multiplication without inserting that conclusion by definition.
\end{remark}

\section{Obstacle 1: Hilbert-Space Completions on Both Sides}
\label{v4-arith:w9-r4:obstacle-hilbert}

The first obstruction is the choice of Hilbert-space completions on
which the Mellin-Hadamard symbol lands as a closable operator. Both
sides have canonical candidates; the candidates are mutually
incompatible.

\subsection{Source-Side Completion}
\label{v4-arith:w9-r4:obstacle-hilbert:source}

\begin{definition}[source-side candidate Hilbert space]
\label{v4-arith:w9-r4:def:source-completion}
\ClaimStatusDefinitional. In the weighted coordinates, the source completion is represented by
\(\ell^{2}(\mathbb{N}^{\times},dm/m)\), isometric to
\(\mathcal H_{\mathrm{BC}}\) through \(\delta_m\mapsto m^{-1/2}v_m\) as in
\Cref{v4-arith:w9-r4:rem:bc-structure}. \(N_{\mathrm{occ}}\) acts as
multiplication by \(\log m\) and is essentially self-adjoint on the
finite-support core
(\Cref{v4-arith:w5-b:occupation-comparison}).
\end{definition}

\begin{lemma}[Mellin normalization and occupation distributions]
\label{v4-arith:w9-r4:lem:no-mellin-on-HBC}
\ClaimStatusProvedHere. For real $c$, the normalized transform
\[
 f\longmapsto (2\pi)^{-1/2}\int_0^\infty f(x)x^{c+i\tau-1}\,dx
\]
extends from compact smooth tests to a unitary
$L^2(\mathbb R_{>0},x^{2c-1}dx)\to L^2(\mathbb R,d\tau)$.
Thus the line $c=1/2$ uses $dx$, while $c=0$ uses $dx/x$.
The occupation map
\[
 Tv_m=m^{-1}\delta_{m^{-1}},\qquad
 \mathcal M(Tv_m)(s)=m^{-s},\qquad
 (-\log x)T=TN_{\mathrm{occ}}
\]
is an injective map into distributions, with additive convention
$\langle\delta_a,\phi\rangle=\phi(a)$.
Every nonzero finite occupation combination has a Mellin transform
outside $L^2$ on every vertical line.
\end{lemma}
\begin{proof}
Under $t=\log x$, $g(t)=e^{ct}f(e^t)$ identifies the source
norm with $L^2(dt)$. The transform is the Fourier unitary with
kernel $(2\pi)^{-1/2}e^{i\tau t}$, as proved in
Lemma~\ref{v4-arith:mo:mellin-unitary}. Evaluation on a point mass
gives the two occupation identities. On a fixed vertical line,
a finite occupation combination is a nonzero exponential polynomial
in $\tau$. It is not square-integrable by
Lemma~\ref{v4-arith:mo:exponential-polynomials}.
\end{proof}

\begin{remark}[the occupation carrier]
\label{v4-arith:w9-r4:rem:lattice-ill-suited}
The distributional map $T$ is defined on occupation $\ell^2$ as
a locally finite sum. Its Dirichlet series converges on
$\operatorname{Re}s>1/2$, with every spectral derivative locally
uniformly convergent. Boundary continuation requires additional
data. These are the precise domain assertions of
Proposition~\ref{v4-arith:mo:dirichlet-domain}.
An additive Dirac distribution is not a vector of continuous
$L^2(dx)$ or $L^2(dx/x)$. Thus $T$ is not an inclusion into
either Mellin Hilbert space.
\end{remark}

\subsection{Target-Side Completion}
\label{v4-arith:w9-r4:obstacle-hilbert:target}

\begin{definition}[target-side candidates]
\label{v4-arith:w9-r4:def:target-completion}
\ClaimStatusDefinitional. Three candidate Hilbert spaces for
\(\mathcal{H}_{\xi}\) in the non-formal sense:
\begin{enumerate}
\item \(\mathcal{H}_{\xi}^{\mathrm{idele}}=L^{2}(C_{\mathbb{Q}},d^{\times}x)\)
the idele class Hilbert space of
\Cref{v4-arith:w7-j:def:H-idele};
\item $K_Z=\ell^2(Z)$ with the signed form $\langle u,Jv\rangle$
and test map $E$ of Theorem~\ref{v4-arith:sc:zero-representation}.
This is a constructed zero-evaluation carrier. It has no specified
identification with an adele-class quotient. Such an identification
requires a separate measure or representation and a domain-preserving
map, as in \Cref{v4-arith:w7-j:def:connes-ac-space};
\item \(\mathcal{H}_{\xi}^{\mathrm{Petersson}}\) the Petersson
completion of the sph-finite cuspidal core
(\Cref{v4-arith:thm:P3-resolution}).
\end{enumerate}
\end{definition}

\begin{lemma}[Mellin transform on the idele norm factor]
\label{v4-arith:w9-r4:lem:mellin-on-idele}
\ClaimStatusProvedHereConditional. Fix a product decomposition
$C_{\mathbb Q}\simeq C_{\mathbb Q}^1\times\mathbb R_{>0}$
and the product Haar measure $dk\,dx/x$, with $dk$ of mass one.
The transform on the norm factor is
\[
 (\mathcal U_0f)(k,\tau)=\frac1{\sqrt{2\pi}}
             \int_0^\infty f(k,x)x^{i\tau}\,\frac{dx}{x}.
\]
It is unitary onto $L^2(C_{\mathbb Q}^1\times\mathbb R,dk\,d\tau)$.
The normalized generator is $i x\partial_x$ on its tensor core.
Writing the same transform on $\operatorname{Re}s=1/2$ requires
replacing $f(k,x)$ by $x^{-1/2}f(k,x)$ and using $dx$ in that
weighted source coordinate.
\end{lemma}
\begin{proof}
The $c=0$ case of Lemma~\ref{v4-arith:mo:mellin-unitary}, tensored
with the identity on $L^2(C_{\mathbb Q}^1)$, proves unitarity.
Substitution into the $c=1/2$ Mellin integral proves the last formula.
The closed generator and its domain follow from multiplication
by $\tau$ under this unitary.
\end{proof}

\begin{remark}[operator comparison between the completions]
\label{v4-arith:w9-r4:rem:completion-mismatch}
The atomic and continuous Hilbert spaces are separable and have
infinite dimension, so they admit abstract unitary isomorphisms.
The specified dynamics impose a stronger condition: an occupation
eigenvector cannot map to a nonzero continuous Mellin eigenvector.
\Cref{v4-arith:w9-r4:atomic-continuous-obstruction} proves
the bounded group-intertwiner obstruction. The explicit distribution
map preserves logarithmic multiplication, while normalized dilation
induces the full residue-jet operator of
Proposition~\ref{v4-arith:mo:jet-comparison}. Their compositions are distinct.
\end{remark}

\section{Obstacle 2: Kernel of \(A\) as Bounded/Unbounded Operator}
\label{v4-arith:w9-r4:obstacle-kernel}

Even accepting a compatible Hilbert-space completion, \(A\) must have
an explicit kernel making it closable. The Hadamard inversion step is
the bottleneck.

\subsection{The Hadamard Inversion Kernel}
\label{v4-arith:w9-r4:obstacle-kernel:had}

\begin{definition}[residue contours and full principal parts]
\label{v4-arith:w9-r4:def:formal-had-kernel}
\ClaimStatusDefinitional. Choose a positively oriented small circle
$C_\rho$ around a zero $\rho$, enclosing no other zero. For an
entire numerator $F$ and a zero of order $r_\rho$, put
\[
 R_{\rho,k}F=\frac1{2\pi i}\int_{C_\rho}
       \frac{(s-\rho)^kF(s)}{\xi_{\mathbb R}(s)}\,ds,
 \qquad 0\leq k<r_\rho.
\]
These are the full principal-part coordinates. Inserting an
additional factor $(s-\rho)^{-1}$ would instead extract the
constant Laurent coefficient of $F/\xi_{\mathbb R}$.
\end{definition}

\begin{lemma}[residue evaluation in the Mellin norm]
\label{v4-arith:w9-r4:lem:had-not-HS}
\ClaimStatusProvedHere. For every real $c$ and every nonempty finite
zero set, full residue evaluation on compact smooth tests has
dense graph in
$L^2(\mathbb R_{>0},x^{2c-1}dx)\oplus\mathscr J_Z$.
It has no closable extension agreeing on these tests. In particular
it is neither bounded nor Hilbert--Schmidt in this source norm.
\end{lemma}
\begin{proof}
This is Theorem~\ref{v4-arith:mo:finite-graph}. For a single simple zero
$c+i\gamma$, its decisive sequence is
$f_L(x)=x^{-c}L^{-1}\psi((\log x)/L)x^{-i\gamma}$,
where $\psi$ is a compact smooth function of integral one.
Its squared source norm is $L^{-1}\|\psi\|_2^2$, while its
Mellin value at that zero remains one. The full finite-jet proof
uses independent exponential-polynomial kernels and requires
neither simplicity nor placement on the line.
\end{proof}

\begin{remark}[the target norm matters]
\label{v4-arith:w9-r4:rem:unboundedness-intrinsic}
Finite-coordinate nonclosability alone does not imply
nonclosability of an arbitrary countable vector-valued map.
Theorem~\ref{v4-arith:mo:countable-completion} and Theorem~\ref{v4-arith:mo:all-jets}
give countable hypotheses and explicit weights for which the
whole graph is dense. Their quotient evaluation norms instead
give well-defined Hilbert spaces and closed dilation operators.
\end{remark}

\subsection{Occupation domains and analytic continuation}
\label{v4-arith:w9-r4:obstacle-kernel:domain}

\begin{proposition}[occupation domains for residue evaluation]
\label{v4-arith:w9-r4:prop:candidate-domains}
\ClaimStatusProvedHere. The finite occupation core maps under $T$
to compact distributions with entire Dirichlet-polynomial Mellin
transforms. If $|a_m|=O(m^{-N})$ for every $N$, then
$F_a(s)=\sum_m a_m m^{-s}$ is entire and bounded on each closed
vertical strip. This coefficient condition does not imply vertical
decay. Given positive residue weights $w_\rho$, the precise
scalar Hilbert domain is
\[
 \{a\in\ell^2:F_a\text{ extends holomorphically near every selected zero},
       \ \sum_\rho w_\rho|R_{\rho,0}F_a|^2<\infty\}.
\]
The source norm and convergence convention for $F_a$ are specified
before taking this domain. On occupation $\ell^2$, its initial
holomorphic domain is $\operatorname{Re}s>1/2$.
\end{proposition}
\begin{proof}
On a fixed strip, choose $N$ exceeding the absolute values of
its real endpoints by two. Absolute convergence bounds the series
uniformly there and gives entire continuation. The vector $v_1$
has $F_{v_1}=1$ and satisfies every coefficient decay condition,
but its transform does not decrease on a vertical line.
Square summability is the definition of membership in the chosen
residue target. It does not follow from coefficient decay without
bounds on the inverse Laurent coefficients at the divisor.
The half-plane assertion is Proposition~\ref{v4-arith:mo:dirichlet-domain}.
\end{proof}

\begin{remark}[closure and the occupation identity]
\label{v4-arith:w9-r4:rem:no-domain-gives-SA}
The finite-core residue identity already fails algebraically:
Proposition~\ref{v4-arith:mo:occupation-defect} computes its defect, and
Proposition~\ref{v4-arith:mo:occupation-jet-graph} shows that no finite full-jet
quotient carries the induced occupation operator.
Under summable positive simple-zero weights,
Proposition~\ref{v4-arith:mo:occupation-positive-obstruction} excludes symmetry
on the finite occupation core even when all selected zeros lie
on the critical line. A different dense domain and a different
comparison map require their own operator construction.
For a closed densely defined map $A$, its natural operator $A^*A$
is positive self-adjoint and its polar decomposition has a partial
isometry. The arithmetic requirement is the construction of such
an $A$ with the stipulated domains, dense range, injectivity,
energy identity, and trace comparison.
\end{remark}

\section{Connes 1999 Partial Approach and Its Limitation}
\label{v4-arith:w9-r4:connes-1999}

Connes' adele class space construction gives an explicit Hilbert space
and scaling generator together with the trace formula. The
Hilbert--P\'olya input is the further assertion that its absorption
trace distribution is the non-trivial zero distribution, with
critical-line support.

\subsection{The Connes Construction}
\label{v4-arith:w9-r4:connes-1999:construction}

\begin{definition}[Connes adele class data]
\label{v4-arith:w9-r4:def:connes-data}
\ClaimStatusDefinitional \emph{(Connes 1999 arXiv:math/9811068 \S IV)}.
The \emph{Connes adele class Hilbert space} is
\(\mathcal{H}_{\mathrm{Connes}}=L^{2}(X_{\mathbb{Q}},dx)\) with
\(X_{\mathbb{Q}}=\mathbb{A}/\mathbb{Q}^{\times}\) the adele class space
and \(dx\) the Haar measure inherited from \(\mathbb{A}\) modulo
\(\mathbb{Q}^{\times}\); the quotient is non-Hausdorff, and the
\(L^{2}\)-structure is realised via the crossed product
\(C^{*}(\mathbb{A}/\mathbb{Q}^{\times})=C_{0}(\mathbb{A})\rtimes\mathbb{Q}^{\times}\).
The \emph{Connes scaling generator} is
\(D_{\mathrm{Connes}}=-i\partial_{\log\lambda}U(\lambda)\) at
\(\lambda=1\), where \(U(\lambda)f(x)=|\lambda|^{-1/2}f(\lambda^{-1}x)\).
On the cuspidal subspace
\(\mathcal{E}\subset\mathcal{H}_{\mathrm{Connes}}\) (orthogonal to the
Tate-regularised constants), \(D_{\mathrm{Connes}}\) is conjectured to
have absorption trace equal to the zero distribution of
\(\xi_{\mathbb R}\). Under zero-placement this trace is supported at
the ordinate multiset \(\{\gamma_{\rho}\}\).
\end{definition}

\begin{theorem}[Connes trace formula equivalence]
\label{v4-arith:w9-r4:thm:connes-trace-eq}
\ClaimStatusProvedElsewhere \emph{(Connes 1999 \S V Thm.~1;
\Cref{v4-arith:tfe:thm:connes-meyer})}. For
\(h\in\mathcal{H}_{\mathrm{test}}\) (the even-Schwartz test class on
\(\mathbb{R}\)), the regularised trace
\(\mathrm{Tr}_{\mathrm{reg}}(h(D_{\mathrm{Connes}})|_{\mathcal{E}})\)
equals the Weil distribution side, conditional on the
absorption-trace identification. With zero-placement, the regularised
trace is \(\sum_{\rho}h(\gamma_{\rho})\).
\end{theorem}

\subsection{The Limitation: Absorption Trace as RH}
\label{v4-arith:w9-r4:connes-1999:limitation}

\begin{theorem}[Connes absorption-trace conjecture]
\label{v4-arith:w9-r4:thm:connes-absorption}
\ClaimStatusConjectured \emph{(Connes 1999 \S IV Thm.~1 conditional
form)}. The scaling generator \(D_{\mathrm{Connes}}\) on
\(\mathcal{E}\subset\mathcal{H}_{\mathrm{Connes}}\) has absorption
trace equal to the non-trivial zero distribution of
\(\xi_{\mathbb R}\); after zero-placement this distribution is
supported at the multiset
\(\{\gamma_{\rho}\mid\rho\in\mathcal{Z}\}\) with multiplicity.
\end{theorem}

\begin{remark}[the source of an absorption comparison]
\label{v4-arith:w9-r4:rem:connes-not-uncond}
An operator on an adelic Hilbert space has its identity map as an
endomorphism of that space. This does not give a map from
$\mathcal H_{\mathrm{BC}}$ into it. An occupation comparison requires
a separately defined map $A_{\mathrm C}$, its source domain, and
$A_{\mathrm C}N_{\mathrm{occ}}=H_{\mathrm C}A_{\mathrm C}$ on a
specified invariant subdomain. An absorption character and this
operator identity are distinct requirements.
\end{remark}

\begin{proposition}[transport to an adelic Hilbert realization]
\label{v4-arith:w9-r4:prop:connes-eq-HPAr}
\ClaimStatusProvedHereConditional. Suppose an adelic Hilbert realization
$(\mathcal H_{\mathrm C},H_{\mathrm C})$ and a Hilbert--P\'olya
realization $(\mathcal H_{\mathrm{HP}},H_{\mathrm{HP}})$ have the
unitary domain comparison $U_{\mathrm C}$ of
Proposition~\ref{v4-arith:mc:transport}.
Then $A_{\mathrm C}=U_{\mathrm C}^{-1}A_{\mathrm{HP}}$ transports
a complete occupation comparison on the same specified source domain.
Trace-class characters and equally normalized determinants transfer
under Proposition~\ref{v4-arith:mc:transport} and
Corollary~\ref{v4-arith:mc:determinant}.
Equality of an absorption character with this operator trace is
an additional arithmetic identity on its stated test class.
\end{proposition}
\begin{proof}
The graph map $1\oplus U_{\mathrm C}^{-1}$ transports the occupation
map. The exact domain conjugacy transports its energy identity and
spectral resolution. The two cited transport results prove the trace
and determinant clauses. They use no identification of a regularized
absorption functional with an ordinary operator trace.
\end{proof}

\begin{remark}[the role of the adelic construction]
\label{v4-arith:w9-r4:rem:adele-transfers}
An independently constructed adelic carrier specifies a new source
for arithmetic geometry. To compare it with the occupation model,
one must prove the core isometry, operator-domain identity, and
character comparison in the preceding proposition. The identity map
on either target space cannot replace those maps.
\end{remark}

\section{Meyer 2005 Partial Approach and Its Limitation}
\label{v4-arith:w9-r4:meyer-2005}

Meyer's global representation supplies a character formula and
spectral multiplicities on a nuclear bornological carrier. Continuous
Mellin scaling supplies a separate Hilbert-space calculation. The
comparison below specifies the maps needed to relate their operators.

\subsection{The Meyer Construction}
\label{v4-arith:w9-r4:meyer-2005:construction}

\begin{definition}[continuous norm test space]
\label{v4-arith:w9-r4:def:meyer-data}
\ClaimStatusDefinitional. In the product Hilbert space of
\Cref{v4-arith:w9-r4:lem:mellin-on-idele}, take the algebraic
tensor core
$\mathcal S_{\mathrm{norm}}=L^2(C_{\mathbb Q}^1)\otimes
C_c^\infty(\mathbb R_{>0})$ and the operator
$D_{\mathrm{norm}}=1\otimes ix\partial_x$ there.
This specifies the continuous scaling calculation.
An identification of a global arithmetic representation with
this Hilbert space requires an additional map.
\end{definition}

\begin{theorem}[continuous Mellin scaling and the global character]
\label{v4-arith:w9-r4:thm:meyer-intertwiner}
\ClaimStatusProvedHere. On the norm factor
$L^2(\mathbb R_{>0},dx/x)$, the unitary $\mathcal U_0$ of
\Cref{v4-arith:w9-r4:lem:mellin-on-idele} conjugates the closure
of $ix\partial_x$ to multiplication by $\tau$.
This operator has continuous spectrum and no eigenvectors.

Meyer's global virtual representation on nuclear bornological
spaces has a different carrier. Its character equals the Weil
distribution, and its spectral multiplicities equal the pole
orders of the completed $L$-function. These are Theorems~5.8
and~5.11 of Meyer (2005), arXiv:math/0311468v3, pp.~43 and~47.
\end{theorem}
\begin{proof}
The operator assertion is the $c=0$ case of
Lemma~\ref{v4-arith:mo:mellin-unitary}. The two assertions about the
global virtual representation are the cited primary theorems.
No identification of that representation with this continuous
Hilbert space is used.
\end{proof}

\subsection{Local multiplicity and a Hilbert comparison}
\label{v4-arith:w9-r4:meyer-2005:limitation}

\begin{remark}[the local quotient map]
\label{v4-arith:w9-r4:rem:meyer-not-A}
At a zero $\rho$ of order $r$, the map
$\mathcal O_\rho/(\xi_{\mathbb R})\to\mathbb C^r$ taking full
principal parts is an isomorphism. Multiplication by $s$ acts
as $\rho I+S$, with $(Sa)_k=a_{k+1}$.
This follows directly from the invertible Taylor-to-residue
matrix in Equation~\eqref{v4-arith:mo:jet-formula}.
The trace on this block equals the trace on $r$ scalar copies
of $\rho$, although their nilpotent actions differ when $r>1$.
Thus a character identity alone does not identify the full
representations or their positive pairings.
\end{remark}

\begin{proposition}[an operator-compatible arithmetic comparison]
\label{v4-arith:w9-r4:thm:meyer-eq-HPAr}
\ClaimStatusProvedHereConditional. Let $\mathscr D_N$ be an
$N_{\mathrm{occ}}$-invariant subspace of its domain. Suppose linear
maps $I:\mathscr D_N\to V$, $U:V\to W$, and
$L:W\to\operatorname{Dom}H$ satisfy
\[
 IN_{\mathrm{occ}}=DI,\qquad UD=KU,\qquad LK=HL
\]
on their stated invariant domains, where $H$ is self-adjoint
on a specified positive Hilbert space. Then $A=LUI$ satisfies
$AN_{\mathrm{occ}}=HA$ on $\mathscr D_N$.
Its closedness, injectivity, dense range, zero-divisor multiplicities,
and arithmetic trace identity require separate hypotheses or proofs.
\end{proposition}
\begin{proof}
Domain compatibility makes the compositions well-defined and gives
$AN_{\mathrm{occ}}=LUIN_{\mathrm{occ}}=LUDI=LKUI=HLUI=HA$.
This computation asserts none of the additional analytic properties.
\end{proof}

\begin{remark}[the positive Hilbert comparison]
\label{v4-arith:w9-r4:rem:meyer-schwartz-only}
The local quotient map preserves the full zero multiplicity.
Theorem~\ref{v4-arith:mo:all-jets} completes these local data in an
explicit positive norm. The resulting dilation is symmetric
exactly for simple zeros on the chosen line.
The construction chooses its norm from the divisor and interpolation
data. An arithmetic realization additionally requires the actual
global pairing and character maps, with their domains and continuity.
Meyer's character theorem and the continuous Mellin unitary do not
construct the occupation map in the preceding proposition.
\end{remark}

\input{chapters/arithmetic/87b_hilbert_comparisons}

\section{Residue Factorization and Arithmetic Comparison}
\label{v4-arith:w9-r4:residual}

A scaling model and an occupation model are related by three
separately specified maps. Their operator identity is proved on the
common invariant domain below; an arithmetic trace requires further data.

\subsection{The Three Comparison Maps}
\label{v4-arith:w9-r4:residual:statement}

\begin{theorem}[conditional passage through two test spaces]
\label{v4-arith:w9-r4:thm:irreducible-residual}
\ClaimStatusProvedHereConditional. Let $\mathcal D_N$ be an
$N_{\mathrm{occ}}$-invariant subspace of its operator domain.
Let $V,W$ be complex vector spaces with endomorphisms $D_V,D_W$.
Suppose linear maps $I:\mathcal D_N\to V$, $U:V\to W$, and
$T:W\to\operatorname{Dom}H_{\mathrm{HP}}$ satisfy
\[
 IN_{\mathrm{occ}}=D_VI,\qquad UD_V=D_WU,\qquad
 TD_W=H_{\mathrm{HP}}T.
\]
Then $A=TUI$ satisfies $AN_{\mathrm{occ}}=H_{\mathrm{HP}}A$ on
$\mathcal D_N$. Closedness, injectivity, dense range, and the
arithmetic character are additional conditions. For the continuous
specialization, fix real $c$ and take $V=\mathscr D$,
$D_V=P_c$, $W=\mathcal M\mathscr D\subset\mathcal O(\mathbb C)$,
$U=(2\pi)^{-1/2}\mathcal M$, and
$(D_WF)(s)=i(c-s)F(s)$. These data construct the middle identity.
\end{theorem}
\begin{proof}
All operators act on the displayed spaces, so
$AN_{\mathrm{occ}}=TUD_VI=TD_WUI=H_{\mathrm{HP}}A$.
The continuous specialization follows from
\eqref{v4-arith:mo:dilation}; $P_c$ preserves $\mathscr D$, so
$D_W$ preserves $W$. It does not change the occupation
operator into that dilation generator.
\end{proof}

For the continuous Hilbert specialization, the first map \(I\)
is not a bounded intertwining embedding of the occupation Hilbert
space into the continuous scaling Hilbert space,
by \Cref{v4-arith:w9-r4:atomic-continuous-obstruction}. Its
possible distributional realization and the last map \(T\) remain
separate constructions. The latter must identify the original zero
divisor, not merely project ordinates to the critical line. These
requirements are in addition to any absorption-trace or positivity
condition imposed on a separately constructed scaling space.

\begin{remark}[finite and archimedean comparison maps]
\label{v4-arith:w9-r4:rem:residual-at-infty}
The occupation shift \(\mu_p\) is a proper isometry; it is not the
unitary multiplication operator \(p^{-it}\) on a Mellin line.
For instance \(\mu_p^*\mu_p=1\) and
\(\mu_p\mu_p^*\ne1\), whereas multiplication by \(p^{-it}\) has both
products equal to one. Any comparison using a unitary dilation must
state the enlarged space and its compression. The local and finite-prime
phase comparisons of \Cref{v4-arith:thermal:prop:two-prime-phase} likewise
precede an adelic identification. Absorption trace remains an
analytic question on its specified scaling carrier; the occupation
construction does not make it the sole comparison required for
arithmetic realization.
\end{remark}

\subsection{The operator and arithmetic comparison data}
\label{v4-arith:w9-r4:residual:table}

The occupation operator is self-adjoint on its maximal weighted
sequence domain. The continuous norm generator is unitarily Mellin
diagonalized on $L^2(dx/x)$. The full residue map instead has the
jet action of Proposition~\ref{v4-arith:mo:jet-comparison}, and its
evaluation completions are those of
Theorem~\ref{v4-arith:mo:all-jets}.

To compare a global arithmetic carrier with a Hilbert realization,
specify a complex test space $V$, a positive Hermitian form $h$ on it,
and a linear map $J:V\to\mathcal H$ with dense range and
$\langle Jv,Jw\rangle=h(v,w)$. If $h$ has a nullspace, the induced
map is defined on its quotient. A scaling operator $D$ on an invariant
subspace $V_0\subseteq V$ must satisfy
$J(V_0)\subseteq\operatorname{Dom}H$ and $JD=HJ$ there.
Graph-core and closure assertions require the corresponding domain
hypotheses in Proposition~\ref{v4-arith:mc:core-unitary}.

For two such Hilbert realizations, a core isometry satisfying those
hypotheses extends to the actual unitary comparison. A closed
occupation map transfers by $A_2=UA_1$ under
Proposition~\ref{v4-arith:mc:transport}. An arithmetic character
requires, in addition, equality with the specified trace functional
on its complete test class. A regularized determinant requires the
initial trace-class domain, branch, and continuation of
Corollary~\ref{v4-arith:mc:determinant}.

Meyer's global character theorem fixes a nuclear bornological
representation. It does not provide $h$, $J$, or the occupation map.
If a transported Weil functional is positive,
Proposition~\ref{v4-arith:mc:positive-form} constructs its cyclic
Hilbert space; comparison with the global representation and its
ordinary trace multiplicities remains additional data.

\subsection{The residue identity and the occupation comparison}
\label{v4-arith:w9-r4:residual:why-irreducible}

\begin{theorem}[the residue identity and its required comparison maps]
\label{v4-arith:w9-r4:thm:irreducibility}
\ClaimStatusProvedHere. The continuous scalar residue identity
requires the stated first-order interpolation condition at every
simple zero. A factorization of an occupation intertwiner through
that residue map additionally requires an invariant occupation
domain and both operator comparisons displayed below.
The explicit distributional embedding constructs the first
occupation-to-distribution map, but its residue composition fails
the occupation energy identity on the finite occupation core.
\end{theorem}

\begin{proof}

The residue test has a specific continuous domain. Put
$\mathscr D=C_c^\infty(\mathbb R_{>0})$, and let $\mathcal Z_1$
be the set of simple non-trivial zeros of $\xi_{\mathbb R}$.
For $f\in\mathscr D$, write $F=\mathcal Mf$ and define
\[
 (A_1f)_\rho=\frac{F(\rho)}{\xi_{\mathbb R}'(\rho)},\qquad
 (M_\gamma a)_\rho=\gamma_\rho a_\rho
 \quad(\rho\in\mathcal Z_1).
\]
Here $A_1$ is the simple-zero restriction of $A^{\mathrm{sym}}$,
with codomain the vector space $\mathbb C^{\mathcal Z_1}$ of all
sequences. No summability norm is imposed. The map $M_\gamma$ is
defined on this entire sequence space, while
$M_{\log}f=-(\log x)f$ preserves $\mathscr D$.
\Cref{v4-arith:w9-r4:lem:formal-intertwine} gives the exact identity
\[
 \bigl(A_1M_{\log}f-M_\gamma A_1f\bigr)_\rho
 =\frac{-F'(\rho)-\gamma_\rho F(\rho)}
        {\xi_{\mathbb R}'(\rho)}.
\]
Consequently the continuous residue relation holds for $f$ exactly
when its numerator vanishes at every $\rho\in\mathcal Z_1$.
Locating the zeros on the critical line does not impose this
condition on $F$. Multiple zeros require the full residue and
jet conditions, not this simple-zero formula.

The stronger atomic target remains a different composition. It asks
for a specified $N_{\mathrm{occ}}$-invariant domain
$\mathscr D_{\mathrm{occ}}\subset\operatorname{Dom}N_{\mathrm{occ}}$
and maps from that domain to a continuous function space and then
to the actual zero-operator carrier. For a factorization through the
present residue test, one would need a linear map
$T:\mathscr D_{\mathrm{occ}}\to\mathscr D$ satisfying
$TN_{\mathrm{occ}}=M_{\log}T$. One would also need
$A_1M_{\log}T=M_\gamma A_1T$. This makes the sequence subspace
$E=A_1T(\mathscr D_{\mathrm{occ}})$ invariant under $M_\gamma$,
since $\mathscr D_{\mathrm{occ}}$ is invariant under $N_{\mathrm{occ}}$.
Finally require a complex-linear map
$L:E\to\operatorname{Dom}H_{\mathrm{HP}}$ satisfying
$LM_\gamma=H_{\mathrm{HP}}L$ on $E$. Only then does
$A=LA_1T$ satisfy
\[
 AN_{\mathrm{occ}}
 =LA_1M_{\log}T=LM_\gamma A_1T=H_{\mathrm{HP}}A.
\]
These equalities specify the domains of every composition. They
supply neither injectivity nor a closed operator realization.
Moreover, if $v_m\in\mathscr D_{\mathrm{occ}}$, the first identity
forces $(-\log x-\log m)Tv_m(x)=0$. A smooth function supported
at the single point $x=m^{-1}$ is zero, so $Tv_m=0$.
Thus this particular compactly supported smooth realization cannot
be injective on a domain containing an occupation eigenvector.
The distributional map of Proposition~\ref{v4-arith:mo:occupation} constructs
that edge on every finite occupation sequence. Its residue composition
still has the nonzero defect of Proposition~\ref{v4-arith:mo:occupation-defect}.
The continuous dilation comparison instead uses $P_c=i(x\partial_x+c)$
and the full jets of Proposition~\ref{v4-arith:mo:jet-comparison}.
\end{proof}

\begin{remark}[the remaining arithmetic construction]
\label{v4-arith:w9-r4:rem:beilinson-closure}
The logarithmic occupation Hamiltonian and the continuous Mellin
generator have the explicitly different actions proved above.
A complete arithmetic realization must construct the global pairing,
the occupation map, and the character and determinant comparisons.
Theorems~\ref{v4-arith:mo:all-jets} and
\ref{v4-arith:w9-r4:thm:irreducibility} give the residue constructions
and their limits. They do not assert that every further arithmetic
comparison reduces to one positivity or scattering statement.
\end{remark}

\section{Placement in Triple-Identification Framework}
\label{v4-arith:w9-r4:placement}

The spectral-flow problem of \Cref{v4-arith:w9-r4:def:spectral-flow}
has proposed comparisons with the constructions in
Chapters~\ref{v4-arith:w8-a:deninger-chiral-homology},
\ref{v4-arith:w7-j:ATP-adelic-scattering}, and
\ref{v4-arith:3d-acs-E1bar}; the following statements specify their Hilbert maps and
operator domains.

\subsection{Spectral Flow as Deninger Spectral-Flow Theorem}
\label{v4-arith:w9-r4:placement:deninger}

\begin{proposition}[transport to a middle-cohomology realization]
\label{v4-arith:w9-r4:prop:spectral-flow-eq-deninger}
\ClaimStatusProvedHereConditional. Suppose a positive Hilbert completion
$\mathcal H_{\mathrm{Den}}$ of the specified primary middle cohomology
and a self-adjoint $H_{\mathrm{Den}}$ have been constructed.
Fix a unitary $U_{\mathrm{Den}}:\mathcal H_{\mathrm{Den}}
\to\mathcal H_{\mathrm{HP}}$ carrying their complete operator domains
and satisfying $U_{\mathrm{Den}}H_{\mathrm{Den}}
=H_{\mathrm{HP}}U_{\mathrm{Den}}$.
Then the occupation map on this cohomological realization is
$A_{\mathrm{Den}}=U_{\mathrm{Den}}^{-1}A_{\mathrm{HP}}$.
It has the same source domain, closedness, injectivity, dense-range,
and intertwining properties. The trace and determinant comparisons
require the hypotheses of Proposition~\ref{v4-arith:mc:transport}
and Corollary~\ref{v4-arith:mc:determinant}.
\end{proposition}
\begin{proof}
Apply Proposition~\ref{v4-arith:mc:transport} to
$U_{\mathrm{Den}}^{-1}$. A core isometry can construct this unitary
only when it compares the actual positive pairings and the graph
cores, as in Proposition~\ref{v4-arith:mc:core-unitary}.
A chain or cohomology identification alone has no such Hilbert
norm or adjoint-domain consequence.
\end{proof}

\begin{remark}[cohomology and Hilbert structure]
\label{v4-arith:w9-r4:rem:deninger-renames}
The middle-cohomology object, its positive pairing, and its completed
scaling action are separate data. The preceding proposition states
the exact map that would identify its operator problem with the
specified Hilbert--P\'olya problem. Existence of that map is not
a consequence of renaming the shifted generator.
\end{remark}

\subsection{Spectral Flow as A-TP Adelic-Scattering}
\label{v4-arith:w9-r4:placement:atp}

\begin{proposition}[scattering and arithmetic comparisons]
\label{v4-arith:w9-r4:prop:spectral-flow-eq-atp}
\ClaimStatusProvedHereConditional. Suppose a scaling realization
$(\mathcal H_a,H_a)$ and a zero-operator realization
$(\mathcal H_\xi,H_\xi)$ have the unitary domain comparison of
Proposition~\ref{v4-arith:mc:transport}.
Then a closed occupation intertwiner, its injectivity and dense range,
and its trace-class character identities transfer between them.
Their equally normalized regularized determinants transfer under
Corollary~\ref{v4-arith:mc:determinant}.

Separately, let $\Omega_\pm:\mathcal K_0\to\mathcal K_a$ be the
isometric wave maps of a specified scattering pair. Their scattering
operator is unitary exactly when their ranges agree.
For an arithmetic test algebra $\mathscr A=C_c^\infty(\mathbb R)$,
a map $C:\mathscr A\to\mathcal K_a$ with dense range and
\[
 W(f*g^*)=\langle Cf,Cg\rangle_{\mathcal K_a}
\]
identifies its positive convolution completion with $\mathcal K_a$.
This pairing identity must hold on the full test algebra, including
any complement of a specified Tate--Mellin test subspace.
\end{proposition}
\begin{proof}
The first assertions are Proposition~\ref{v4-arith:mc:transport}
and Corollary~\ref{v4-arith:mc:determinant}.
Proposition~\ref{v4-arith:mc:scattering} proves the scattering
range criterion. The displayed pairing identity gives positivity
and the completed isometry of
Proposition~\ref{v4-arith:mc:positive-form}.
No implication between the range criterion and the arithmetic
pairing identity is used. For a continuous arithmetic form, equality
on a subspace extends to a larger test space only when that subspace
is dense in the stated test topology and both forms are continuous.
Without these hypotheses, the identity on the larger space is
an additional condition.
\end{proof}

\begin{remark}[the scattering condition does not specify the divisor]
\label{v4-arith:w9-r4:rem:scattering-transfers}
The equality of wave-map ranges is a statement about a specified
scattering pair. The one-dimensional example after
Proposition~\ref{v4-arith:mc:scattering} has unitary scattering and
the wrong xi divisor. Thus an arithmetic equivalence also requires
the carrier isometry, the full pairing identity, and the character
and determinant comparisons of the preceding proposition.
The continuous Mellin unitary supplies none of these arithmetic maps.
\end{remark}

\subsection{Spectral Flow as 3D ACS Bulk-Boundary Correspondence}
\label{v4-arith:w9-r4:placement:acs}

\begin{proposition}[bulk and occupation operator types]
\label{v4-arith:w9-r4:prop:spectral-flow-eq-acs}
\ClaimStatusProvedHereConditional. Suppose a positive bulk Hilbert space
$\mathcal H_{\mathrm B}$ and self-adjoint operator $H_{\mathrm B}$
are provided by a specified arithmetic Chern--Simons construction.
Suppose also that $U_{\mathrm B}:\mathcal H_{\mathrm B}\to
\mathcal H_{\mathrm{HP}}$ is the unitary domain comparison of
Proposition~\ref{v4-arith:mc:transport}.
Then $A_{\mathrm B}=U_{\mathrm B}^{-1}A_{\mathrm{HP}}$ is the
corresponding occupation map, with the transported properties proved
there. A parameter derivative of bulk sections does not by itself
define this map on occupation states or an operator on a single fiber.
\end{proposition}
\begin{proof}
The transport assertions follow from the explicit formula and
Proposition~\ref{v4-arith:mc:transport}.
For the derivative assertion, trivialize a nonzero Hilbert bundle
near a parameter $s_0$. Its $C^1$ sections include
$a(s)=(s-s_0)v$ and the zero section, for $v\ne0$.
Both have value zero at $s_0$, but their derivatives at $s_0$
are $v$ and zero. Thus evaluation of the section derivative does
not factor through the value in the fiber.
A connection or a common-domain family of operators can specify
a derivative on sections; a further map from occupation states
is still required for the displayed occupation identity.
\end{proof}

\begin{remark}[partition functions and bulk comparisons]
\label{v4-arith:w9-r4:rem:acs-wilson}
An identity between a bulk partition function and $\xi_{\mathbb R}$
is a scalar identity. The bulk Hilbert pairing, scaling generator,
occupation map, and ordinary or regularized character require
their own constructions. The arithmetic object
$\mathcal V^{\mathrm{prim}}$ is a Hochschild trace class; its Hilbert
realization is additional data. Under the unitary and trace hypotheses
above, the complete operator realization transfers by the proved maps.
\end{remark}

\subsection{The Unified Residual}
\label{v4-arith:w9-r4:placement:unified}

\begin{theorem}[comparison of specified arithmetic realizations]
\label{v4-arith:w9-r4:thm:unified-residual}
\ClaimStatusProvedHereConditional. Suppose Hilbert realizations
$(\mathcal H_j,H_j)$ have been constructed for the Hilbert--P\'olya,
adelic absorption, global Meyer, and middle-cohomology descriptions,
with self-adjoint $H_j$.
Fix graph cores $\mathcal C_j$ and bijective core isometries
$U_{j,0}:\mathcal C_j\to\mathcal C_{\mathrm{HP}}$ satisfying
$U_{j,0}H_j=H_{\mathrm{HP}}U_{j,0}$ on those cores.
Assume the invariant-core hypotheses of
Proposition~\ref{v4-arith:mc:core-unitary}, or specify the equivalent
unitary comparisons on the complete operator domains.
Then these maps extend uniquely to unitaries $U_j$ and identify
the following operator properties in all four realizations:
\begin{enumerate}
\item a closed injective occupation map on the same specified dense
domain, with dense range and the occupation identity on its
specified $N_{\mathrm{occ}}$-invariant subdomain;
\item the full shifted spectral divisor, including eigenspace
multiplicities;
\item trace-class characters on the same bounded Borel test class;
\item the equally normalized determinant, when the branch and
continuation hypotheses of Corollary~\ref{v4-arith:mc:determinant} hold.
\end{enumerate}
The corresponding occupation maps are explicitly
$A_j=U_j^{-1}A_{\mathrm{HP}}$.

For the fifth description, arithmetic Tate positivity, suppose a
specified continuous Weil functional satisfies the positivity
hypothesis of Proposition~\ref{v4-arith:mc:positive-form}.
That proposition constructs its cyclic Hilbert space. Its comparison
with one of the preceding realizations requires a dense pairing map
$C$ and the associated scaling-domain and character comparisons.
Under those additional data, the same transport assertions apply.
\end{theorem}
\begin{proof}
Complete the core isometries by Proposition~\ref{v4-arith:mc:core-unitary}.
Their graph identities give equality on the self-adjoint domains.
Proposition~\ref{v4-arith:mc:transport} applied to $U_j^{-1}$ proves
the first three claims and the formula for $A_j$.
Corollary~\ref{v4-arith:mc:determinant} proves the fourth.
For the positivity description, the map $[f]\mapsto Cf$ extends
to the unitary in Proposition~\ref{v4-arith:mc:positive-form}.
The explicitly required operator and character identities then
permit the same transport argument.
\end{proof}

\begin{remark}[critical-line placement and the comparison problem]
\label{v4-arith:w9-r4:rem:honest-summary}
If one self-adjoint $H_j$ has shifted spectral divisor equal to
the entire zero divisor of $\xi_{\mathbb R}$, every such zero has
real part $1/2$. This follows because the eigenvalues of $H_j$
are real. It proves the Riemann Hypothesis from that realization.
The converse does not supply the core isometries, the arithmetic
pairing, or a closed occupation map. In particular, a positive
cyclic form alone does not recover ordinary trace multiplicities,
as the finite example after Proposition~\ref{v4-arith:mc:positive-form}
shows. Proposition~\ref{v4-arith:mo:vacuum} excludes an injective
literal occupation map on any domain containing $v_1$.
These are additional construction conditions in the original
spectral-flow question.
\end{remark}

\section{Polyakov--Dirac Normal Form for the Residual}
\label{v4-arith:w9-r4:polyakov-dirac-normal-form}

Under the stated primary Heisenberg hypotheses, the Polyakov
rank-residue calculation fixes the scalar $r_\zeta=2$.
For a specified symmetric core, the Dirac double tests essential
self-adjointness of that same operator. The arithmetic determinant,
positive Hilbert pairing, occupation map, and character comparison
remain independent hypotheses. The following construction keeps
these data separate and proves the fixed-core equivalence.

\subsection{Polyakov Rigidity of the Rank Residue}
\label{v4-arith:w9-r4:pdnf:polyakov}

\begin{proposition}[Polyakov rank-residue rigidity]
\label{v4-arith:w9-r4:prop:polyakov-scalar-rigidity}
\ClaimStatusProvedHere. In the primary Heisenberg domain of
Chapters~\ref{v4-arith:polyakov} and
\ref{v4-arith:w9-r3:chapter}, any spectral-flow closure compatible
with
\begin{enumerate}
\item the rank-\(2\) per-prime Sugawara Heisenberg stress tensor,
\item the global Fock partition
\(Z(s)=\prod_{p}\prod_{n\geq 1}(1-p^{-sn})^{-2}\), and
\item the archimedean Polyakov anomaly of the rank-\(2\) free-field
completion,
\end{enumerate}
has determinant rank residue
\begin{equation}
\label{v4-arith:w9-r4:eq:polyakov-rigid-c}
r_{\zeta}\;=\;2.
\end{equation}
This determines the scalar under the stated primary Heisenberg
hypotheses. The Hilbert pairing, operator domains, and arithmetic
trace comparisons remain the separate data of
Theorem~\ref{v4-arith:w9-r4:thm:unified-residual}.
\end{proposition}

\begin{proof}
\Cref{v4-arith:w9-r3:thm:main} proves that the unrenormalised conformal scalar
\(\sum_{p}c^{(p)}\) diverges, that the negative global modes do not act
on the restricted vacuum, and that the global Fock-partition residue
gives \(r_{\zeta}=2\). The same chapter proves the Polyakov
rank comparison
\Cref{v4-arith:w9-r3:thm:polyakov-identity}, identifying this
finite-prime residue with the rank-\(2\) archimedean conformal
anomaly. These inputs leave no determinant scalar parameter: changing
\(r_{\zeta}\) would either abandon the Euler-product
partition, abandon the Sugawara rank, or abandon the archimedean
anomaly match. None of those changes is a spectral-flow closure
inside the stated primary Heisenberg domain. This scalar argument
contains no Hilbert isometry, graph-core theorem, or arithmetic
character map. Those comparisons retain the independent hypotheses
of Theorem~\ref{v4-arith:w9-r4:thm:unified-residual}.
\end{proof}

\begin{remark}[what Polyakov does and does not buy]
\label{v4-arith:w9-r4:rem:polyakov-no-rh}
Polyakov regularisation is decisive for the determinant rank,
not for RH. It supplies the finite scalar \(r_{\zeta}=2\) and
rules out the divergent prime-zeta prescription; it does not
produce the Hilbert space on which the zero ordinates are a
self-adjoint spectrum.
\end{remark}

\subsection{Dirac Closure of the Remaining Operator}
\label{v4-arith:w9-r4:pdnf:dirac}

\begin{definition}[a specified symmetric core]
\label{v4-arith:w9-r4:def:residual-ordinate-core}
\ClaimStatusDefinitional. Fix a complex positive Hilbert space
$\mathcal H_{\mathrm{res}}$, a dense linear subspace
$\mathcal C_{\mathrm{res}}$, and a symmetric operator
$H_{\mathrm{res}}^{\mathrm{sym}}:\mathcal C_{\mathrm{res}}
\to\mathcal H_{\mathrm{res}}$.
These are the data for the Dirac construction below.

A continuous Mellin chart for these data requires a specified
unitary $B:\mathcal H_{\mathrm{res}}\to\mathscr H_c$ with
$B\mathcal C_{\mathrm{res}}\subseteq\operatorname{Dom}P_c$ and
$BH_{\mathrm{res}}^{\mathrm{sym}}=P_cB$ on the core.
Then $\mathcal U_cB$ represents the core operator by multiplication
by the real frequency. One may construct $B$ by completing a dense
core isometry under Proposition~\ref{v4-arith:mc:core-unitary}.
Neither the continuous Mellin unitary nor the global Meyer character
theorem supplies that isometry from an arithmetic quotient.

A residue boundary condition is a separate map on a specified
test topology. Its full local action is the jet operator of
Proposition~\ref{v4-arith:mo:jet-comparison}. It is not a restriction
of ordinary $L^2$ point evaluation. If all full jets are retained,
positive symmetry requires simple zeros on the chosen line by
Theorem~\ref{v4-arith:mo:finite-symmetry}.
\end{definition}

\begin{definition}[Dirac double of the residual]
\label{v4-arith:w9-r4:def:dirac-double}
\ClaimStatusDefinitional. The \emph{Dirac double} of the residual
ordinate core is the odd symmetric operator
\begin{equation}
\label{v4-arith:w9-r4:eq:dirac-double}
\mathscr{D}_{\mathrm{res}}^{0}
\;:=\;
\begin{pmatrix}
0 & H_{\mathrm{res}}^{\mathrm{sym}}\\
H_{\mathrm{res}}^{\mathrm{sym}} & 0
\end{pmatrix}
\quad\text{on}\quad
\mathcal{C}_{\mathrm{res}}\oplus\mathcal{C}_{\mathrm{res}},
\end{equation}
with grading
\(\Gamma=\operatorname{diag}(1,-1)\). A \emph{Dirac closure} is a
self-adjoint closure
\(\mathscr{D}_{\mathrm{res}}=\overline{\mathscr{D}_{\mathrm{res}}^{0}}\)
whose spectral measure satisfies the regularised determinant
identity
\begin{equation}
\label{v4-arith:w9-r4:eq:dirac-det}
\det\nolimits_{\infty}\bigl(s\cdot\mathrm{id}
-({\tfrac12}\mathrm{id}+iH_{\mathrm{res}})\bigr)
\;=\;\xi_{\mathbb{R}}(s),
\end{equation}
where \(H_{\mathrm{res}}\) is the off-diagonal component recovered
from \(\mathscr{D}_{\mathrm{res}}\) by the grading.
\end{definition}

\begin{lemma}[Dirac double detects self-adjointness]
\label{v4-arith:w9-r4:lem:dirac-double-detects-sa}
\ClaimStatusProvedHere. Let \(H\) be a densely defined symmetric
operator on a Hilbert space \(\mathcal{H}\) and
\(\mathscr{D}_{H}^{0}=
\left(\begin{smallmatrix}0&H\\H&0\end{smallmatrix}\right)\)
on \(\mathcal{D}(H)\oplus\mathcal{D}(H)\). Then
\(\mathscr{D}_{H}^{0}\) is essentially self-adjoint if and only if
\(H\) is essentially self-adjoint. When this holds,
\[
\overline{\mathscr{D}_{H}^{0}}
\;=\;
\begin{pmatrix}
0 & \overline{H}\\
\overline{H} & 0
\end{pmatrix},
\qquad
\overline{\mathscr{D}_{H}^{0}}^{\,2}
\;=\;
\begin{pmatrix}
\overline{H}^{\,2} & 0\\
0 & \overline{H}^{\,2}
\end{pmatrix}.
\]
\end{lemma}

\begin{proof}
For a symmetric \(H\),
\((\mathscr{D}_{H}^{0})^{*}=
\left(\begin{smallmatrix}0&H^{*}\\H^{*}&0\end{smallmatrix}\right)\) on
\(\mathcal{D}(H^{*})\oplus\mathcal{D}(H^{*})\). Hence the closure of
\(\mathscr{D}_{H}^{0}\) is self-adjoint exactly when the closure of
\(H\) has adjoint domain equal to its own domain, i.e.\ exactly when
\(H\) is essentially self-adjoint. The displayed square is the usual
block-matrix product on
$\operatorname{Dom}(\overline H^2)\oplus\operatorname{Dom}(\overline H^2)$,
as given by the spectral theorem for the self-adjoint operator
\(\overline{H}\) (Reed--Simon I, Ch.~VIII).
\end{proof}

\begin{theorem}[Dirac normal form on a specified graph core]
\label{v4-arith:w9-r4:thm:polyakov-dirac-normal-form}
\ClaimStatusProvedHereConditional. Fix the positive Hilbert space
and symmetric operator $T_0$ on $\mathcal C_{\mathrm{res}}$ in
Definition~\ref{v4-arith:w9-r4:def:residual-ordinate-core}.
Use one fixed determinant regularization and polar-factor convention.
The following statements are equivalent:
\begin{enumerate}
\item $T_0$ is essentially self-adjoint and the determinant of
$1/2+i\overline T_0$ exists in that convention and equals
$\xi_{\mathbb R}$;
\item the Dirac double of this same $T_0$ has a self-adjoint closure
and its recovered off-diagonal operator has that determinant.
\end{enumerate}
Suppose further that a unitary $U:\mathcal H_{\mathrm{res}}
\to\mathcal H_{\mathrm{HP}}$ maps $\mathcal C_{\mathrm{res}}$ onto
a graph core of a self-adjoint $H_{\mathrm{HP}}$ and satisfies
$UT_0=H_{\mathrm{HP}}U$ there. Then the two statements are equivalent
to the same determinant identity for $H_{\mathrm{HP}}$.
An arithmetic absorption-character identity transfers only under
the trace hypotheses of Proposition~\ref{v4-arith:mc:transport}
and the identification of that trace with the specified arithmetic
character.
\end{theorem}
\begin{proof}
The fixed unitary matrix
$2^{-1/2}\left(\begin{smallmatrix}1&1\\1&-1\end{smallmatrix}\right)$
conjugates the Dirac double on its exact product domain to
$T_0\oplus(-T_0)$. Closure and adjoint preserve this unitary
conjugacy. Thus essential self-adjointness is equivalent on the two
sides, and the recovered component is exactly $\overline T_0$.
The determinant clauses are consequently the same clause about
the same closed operator.
For the additional comparison, close the graph identity and use
the assumed graph core to obtain
$U\overline T_0U^{-1}=H_{\mathrm{HP}}$.
Corollary~\ref{v4-arith:mc:determinant} transports the determinant
when its analytic hypotheses hold. Trace invariance then gives
only the stated character transfer. The scalar normalization does
not construct $U$ or prove that an arbitrary dense subspace is a
graph core.
\end{proof}

\begin{remark}[graph completion and the arithmetic determinant]
\label{v4-arith:w9-r4:rem:single-estimate-left}
On the specified symmetric core, the graph norm is
\[
 \|u\|_{\mathscr D}^{2}=\|u\|^{2}
       +\|H_{\mathrm{res}}^{\mathrm{sym}}u\|^{2}.
\]
Completing this graph gives a closed operator. Essential
self-adjointness further requires equality with its adjoint domain.
After a closed Tate-polar operator quotient has been constructed,
the same distinction applies to the quotient operator. Its determinant
identity \eqref{v4-arith:w9-r4:eq:dirac-det}, the occupation map,
and the arithmetic character comparison remain separate conditions.
The fixed-core theorem identifies the Dirac and symmetric-operator
conditions; it does not identify all these arithmetic constructions
with graph completeness.
\end{remark}

\begin{definition}[closed Tate-polar operator quotient]
\label{v4-arith:w9-r4:def:tate-polar-operator-quotient}
\ClaimStatusConditional. A \emph{closed Tate-polar operator quotient}
for \(H_{0}=H_{\mathrm{res}}^{\mathrm{sym}}\) is additional data of a
Hilbert quotient
\[
\pi\colon \mathcal{H}_{\mathrm{res}}\longrightarrow
\mathcal{H}_{\mathrm{res}}^{\circ}
\]
with closed finite-rank kernel \(P_{\mathrm{Tate}}\), an induced
dense domain
\[
\mathcal{C}_{\mathrm{res}}^{\circ}:=\pi(\mathcal{C}_{\mathrm{res}}),
\]
and a symmetric operator \(H_{0}^{\circ}\) on
\(\mathcal{C}_{\mathrm{res}}^{\circ}\) satisfying
\[
H_{0}^{\circ}\pi u=\pi H_{0}u,\qquad u\in\mathcal{C}_{\mathrm{res}}.
\]
The quotient is required to be compatible with closure and adjoint
domains:
\[
\mathcal{D}(\overline{H_{0}^{\circ}})
=\pi\mathcal{D}(\overline{H_{0}}),
\qquad
\mathcal{D}((H_{0}^{\circ})^{*})
=\pi\mathcal{D}(H_{0}^{*}),
\]
with the displayed images understood in the corresponding graph
topologies. A unitary Mellin chart on $\mathcal H_{\mathrm{res}}^\circ$
is additional data as in
Definition~\ref{v4-arith:w9-r4:def:meyer-coordinate-core}.
If it is obtained from a chart before quotienting, the target must
also be quotiented by the image of $P_{\mathrm{Tate}}$, and the
induced pairing and operator domains must be compared. An injective
chart cannot itself factor through a quotient with nonzero kernel.

The quotient is not constructed here. Every phrase
``after quotienting the Tate polar directions'' in the residual
operator argument is to be read under this hypothesis.
\end{definition}

\begin{remark}[post-quotient notation]
\label{v4-arith:w9-r4:rem:post-tate-notation}
\ClaimStatusDefinitional. In the boundary and deficiency discussion
below,
\[
T_{0}:=H_{0}^{\circ}\colon
\mathcal{C}_{\mathrm{res}}^{\circ}\subset
\mathcal{H}_{\mathrm{res}}^{\circ}\longrightarrow
\mathcal{H}_{\mathrm{res}}^{\circ}
\]
denotes the post-Tate residual operator supplied by
\Cref{v4-arith:w9-r4:def:tate-polar-operator-quotient}. The
pre-quotient operator \(H_{0}=H_{\mathrm{res}}^{\mathrm{sym}}\) is only
the source operator. Thus every adjoint, deficiency kernel, boundary
quotient, and Mellin-coordinate graph statement below is a statement
about \(T_{0}\). Without the closed Tate-polar quotient, these formulas
are well-formed as conditional operator theory but not as an
unconditional construction.
\end{remark}

\subsection{Boundary Form and Deficiency Kernels}
\label{v4-arith:w9-r4:pdnf:what-remains}

Under the stated primary Heisenberg hypotheses, scalar rigidity
fixes \(r_{\zeta}=2\). For the specified symmetric operator, the
boundary quotient \(\mathcal Q_{\mathrm{res}}\) tests essential
self-adjointness. The arithmetic determinant, pairing, and occupation
map are not determined by this quotient.

\begin{definition}[residual boundary form]
\label{v4-arith:w9-r4:def:residual-boundary-form}
\ClaimStatusDefinitional. Under the post-quotient notation of
\Cref{v4-arith:w9-r4:rem:post-tate-notation}, the adjoint-domain
boundary form is the Hermitian symplectic form
\begin{equation}
\label{v4-arith:w9-r4:eq:boundary-form}
\mathcal{B}_{\mathrm{res}}(u,v)
\;:=\;
\langle T_{0}^{*}u,v\rangle-\langle u,T_{0}^{*}v\rangle,
\qquad
u,v\in\mathcal{D}(T_{0}^{*}).
\end{equation}
The \emph{residual boundary space} is the quotient
\begin{equation}
\label{v4-arith:w9-r4:eq:boundary-space}
\mathcal{Q}_{\mathrm{res}}
\;:=\;
\mathcal{D}(T_{0}^{*})/\mathcal{D}(\overline{T_{0}}),
\end{equation}
with the induced form still denoted \(\mathcal{B}_{\mathrm{res}}\).
\end{definition}

\begin{proposition}[the boundary quotient of a fixed symmetric operator]
\label{v4-arith:w9-r4:prop:polyakov-boundary-target}
\ClaimStatusProvedHere. For the specified densely defined symmetric
operator $T_0$, essential self-adjointness is equivalent to
\begin{equation}
\label{v4-arith:w9-r4:eq:polyakov-boundary-target}
 \mathcal Q_{\mathrm{res}}=0.
\end{equation}
The induced boundary form is nondegenerate on this quotient.
Consequently, a bijection from a specified BV--BFV boundary space
to $\mathcal Q_{\mathrm{res}}$ preserving the complete boundary
pairing would identify vanishing of that pairing with essential
self-adjointness. Cancellation of a variation merely on the original
core proves only its isotropy. Neither conclusion identifies the
arithmetic determinant or constructs an occupation map.
\end{proposition}
\begin{proof}
Since $T_0$ is symmetric, it is closable and
$\overline T_0\subseteq T_0^*$. The quotient vanishes exactly when
their domains agree; their actions then agree, so
$\overline T_0=(\overline T_0)^*$.
If the class of $u\in\operatorname{Dom}T_0^*$ pairs to zero with
every class, then
$\langle T_0^*u,v\rangle=\langle u,T_0^*v\rangle$ for every
$v\in\operatorname{Dom}T_0^*$. The adjoint definition gives
$u\in\operatorname{Dom}(T_0^*)^*=\operatorname{Dom}\overline T_0$.
Thus the quotient form is nondegenerate. A bijection preserving
the pairing transfers its vanishing; isotropy of the original
core alone says nothing about the remaining adjoint-domain classes.

For the missing arithmetic implication, take $T_0=0$ on
$\mathbb C$. Its boundary quotient is zero, while its shifted
operator has eigenvalue $1/2$. Proposition~\ref{v4-arith:mo:vacuum}
proves $\xi_{\mathbb R}(1/2)>1/4$, so this operator has the wrong
zero divisor. Scalar normalization cannot change that divisor.
The occupation and arithmetic character comparisons require
the maps and hypotheses of Proposition~\ref{v4-arith:mc:transport}.
\end{proof}

\begin{definition}[deficiency spaces of the residual]
\label{v4-arith:w9-r4:def:deficiency-spaces}
\ClaimStatusDefinitional. The residual deficiency spaces are
\begin{equation}
\label{v4-arith:w9-r4:eq:deficiency-spaces}
\mathcal{N}_{+}:=\ker(T_{0}^{*}-i),
\qquad
\mathcal{N}_{-}:=\ker(T_{0}^{*}+i).
\end{equation}
\end{definition}

\begin{theorem}[two conditions on the fixed quotient operator]
\label{v4-arith:w9-r4:thm:dirac-two-test}
\ClaimStatusProvedHereConditional. Fix the positive quotient and
symmetric operator $T_0$ above, with the determinant convention of
Theorem~\ref{v4-arith:w9-r4:thm:polyakov-dirac-normal-form}.
Its Dirac closure is equivalent to the following two statements:
\begin{enumerate}
\item \emph{Deficiency vanishing:}
\begin{equation}
\label{v4-arith:w9-r4:eq:deficiency-vanishing}
\mathcal{N}_{+}=0=\mathcal{N}_{-}.
\end{equation}
\item \emph{Determinant identity for the closure:}
\begin{equation}
\label{v4-arith:w9-r4:eq:closure-det}
\det\nolimits_{\infty}
\bigl(s\cdot\mathrm{id}
-({\tfrac12}\mathrm{id}+i\overline{T_{0}})\bigr)
\;=\;
\xi_{\mathbb{R}}(s).
\end{equation}
\end{enumerate}
Equivalently, the Dirac double
\(\mathscr{D}_{\mathrm{res}}^{0}\) has no deficiency states and its
closed off-diagonal component has determinant \(\xi_{\mathbb R}\).
\end{theorem}

\begin{proof}
By von Neumann's deficiency-index criterion, a densely defined
symmetric operator \(T_{0}\) is essentially self-adjoint if and only
if
\(\ker(T_{0}^{*}-i)=\ker(T_{0}^{*}+i)=0\). Thus
\eqref{v4-arith:w9-r4:eq:deficiency-vanishing} is exactly the
essential self-adjointness part of
\Cref{v4-arith:w9-r4:def:dirac-double}. By
\Cref{v4-arith:w9-r4:lem:dirac-double-detects-sa}, the same
condition is equivalent to essential self-adjointness of the Dirac
double \(\mathscr{D}_{\mathrm{res}}^{0}\).

For this specified symmetric core, after essential self-adjointness
the remaining clause of its Dirac closure is the regularised
determinant. Arithmetic carrier and occupation comparisons remain
separate hypotheses. Equation
\eqref{v4-arith:w9-r4:eq:closure-det} is precisely
\eqref{v4-arith:w9-r4:eq:dirac-det} for the closure
\(\overline{T_{0}}\). Combining deficiency vanishing with the
determinant identity is therefore equivalent to the Dirac closure
condition of
\Cref{v4-arith:w9-r4:thm:polyakov-dirac-normal-form}, hence to the
Dirac condition on the specified core. Arithmetic transport further
requires Proposition~\ref{v4-arith:mc:transport}.
\end{proof}

\begin{remark}[the information in deficiency kernels]
\label{v4-arith:w9-r4:rem:smartest-next-calculation}
For the specified symmetric core, the kernels
$\ker(T_0^*\mp i)$ test essential self-adjointness. Their vanishing
gives a self-adjoint closure. A nonzero deficiency vector gives
a nonzero class in $\mathcal Q_{\mathrm{res}}$.
The determinant identity \eqref{v4-arith:w9-r4:eq:closure-det}
and the arithmetic comparisons are independent additional conditions,
as the one-dimensional example in
Proposition~\ref{v4-arith:w9-r4:prop:polyakov-boundary-target} shows.
The scalar identity $r_\zeta=2$ under its stated primary hypotheses
does not determine these kernels or construct the comparison maps.
\end{remark}

\subsection{The Deficiency Calculation}
\label{v4-arith:w9-r4:pdnf:deficiency-calculation}

In specified Mellin coordinates the deficiency kernels are the orthogonal
complements of the ranges of \((M_t \pm i)\) on the residual core,
with the sign convention matching the conjugate-linear second slot in
the inner product. The computation is unconditional given the specified
Mellin chart and the closed Tate-polar operator quotient of
\Cref{v4-arith:w9-r4:def:tate-polar-operator-quotient}; without those
it is conditional operator theory. The obstruction is a graph-core
question for the residual Hadamard domain, not an Euler-product
calculation.

\begin{definition}[Mellin-coordinate residual chart]
\label{v4-arith:w9-r4:def:meyer-coordinate-core}
\ClaimStatusDefinitional. Let
\[
M_{t}\colon \mathcal{D}(M_{t})\subset L^{2}(\mathbb{R},dt)\to
L^{2}(\mathbb{R},dt),\qquad (M_{t}f)(t)=t f(t),
\]
be the self-adjoint multiplication operator. A \emph{Mellin-coordinate
residual chart} is a unitary map, denoted \(U_{\mathrm{res}}\) by
abuse of notation,
\[
U_{\mathrm{res}}\colon\mathcal{H}_{\mathrm{res}}^{\circ}
\longrightarrow L^{2}(\mathbb{R},dt),
\]
from the residual quotient Hilbert space of
\Cref{v4-arith:w9-r4:def:tate-polar-operator-quotient} such that the
Mellin-coordinate residual core is
\begin{equation}
\label{v4-arith:w9-r4:eq:meyer-coordinate-core}
\mathcal{C}_{\mathrm{res}}^{M}
\;:=\;
U_{\mathrm{res}}(\mathcal{C}_{\mathrm{res}}^{\circ})
\;\subset\;
\mathcal{D}(M_{t}),
\end{equation}
and
\[
T_{0}=H_{0}^{\circ}
\quad\text{is represented by}\quad
M_{t}|_{\mathcal{C}_{\mathrm{res}}^{M}} .
\]
This chart is additional Hilbert-space data. The continuous
Mellin theorem provides $\mathcal U_c:\mathscr H_c\to L^2(\mathbb R)$,
so a sufficient construction is a specified unitary
$B:\mathcal H_{\mathrm{res}}^\circ\to\mathscr H_c$ satisfying
$B\mathcal C_{\mathrm{res}}^\circ\subseteq\operatorname{Dom}P_c$
and $BT_0=P_cB$ on that core. Then the chart is the actual composition
$U_{\mathrm{res}}=\mathcal U_cB$.
A dense isometric core map can construct $B$ under
Proposition~\ref{v4-arith:mc:core-unitary}. Once the unitary and
the exact core-domain identity are given, that proposition also
transports adjoint domains and deficiency kernels. The global
Meyer character theorem supplies no such map from the residual quotient.
\end{definition}

\begin{theorem}[deficiency kernels are range defects]
\label{v4-arith:w9-r4:thm:deficiency-range-defect}
\ClaimStatusProvedHereConditional. Assume a Mellin-coordinate residual chart in
the sense of \Cref{v4-arith:w9-r4:def:meyer-coordinate-core}. Then
\begin{equation}
\label{v4-arith:w9-r4:eq:deficiency-range-defect-plus}
U_{\mathrm{res}}\mathcal{N}_{+}
\;=\;
\bigl((M_{t}+i)\mathcal{C}_{\mathrm{res}}^{M}\bigr)^{\perp},
\end{equation}
and
\begin{equation}
\label{v4-arith:w9-r4:eq:deficiency-range-defect-minus}
U_{\mathrm{res}}\mathcal{N}_{-}
\;=\;
\bigl((M_{t}-i)\mathcal{C}_{\mathrm{res}}^{M}\bigr)^{\perp}.
\end{equation}
Consequently \(\mathcal{N}_{+}=0=\mathcal{N}_{-}\) if and only if
both ranges
\[
(M_{t}-i)\mathcal{C}_{\mathrm{res}}^{M},
\qquad
(M_{t}+i)\mathcal{C}_{\mathrm{res}}^{M}
\]
are dense in \(L^{2}(\mathbb{R},dt)\).
\end{theorem}

\begin{proof}
We prove the \(+\) identity; the \(-\) identity is identical with
\(i\) replaced by \(-i\). The inner product is linear in the first
slot and conjugate-linear in the second. By definition,
\(u\in\mathcal{N}_{+}\) means \(u\in\mathcal{D}(T_{0}^{*})\) and
\((T_{0}^{*}-i)u=0\). Equivalently, for every
\(f\in\mathcal{C}_{\mathrm{res}}^{\circ}\),
\[
\langle T_{0}f,u\rangle
\;=\;
\langle f,iu\rangle .
\]
Since the second slot is conjugate-linear,
\(\langle f,iu\rangle=-i\langle f,u\rangle\).
Transport by \(U_{\mathrm{res}}\). Write
\(\tilde u=U_{\mathrm{res}}u\) and
\(\tilde f=U_{\mathrm{res}}f\). Since \(T_{0}\) is represented by
\(M_{t}\) on \(\mathcal{C}_{\mathrm{res}}^{M}\), the adjoint equation
becomes
\[
\langle M_{t}\tilde f,\tilde u\rangle
\;=\;
\langle \tilde f,i\tilde u\rangle ,
\qquad
\tilde f\in\mathcal{C}_{\mathrm{res}}^{M}.
\]
Hence
\[
\langle (M_{t}+i)\tilde f,\tilde u\rangle=0
\qquad
\text{for all }\tilde f\in\mathcal{C}_{\mathrm{res}}^{M},
\]
which says precisely that
\(\tilde u\in((M_{t}+i)\mathcal{C}_{\mathrm{res}}^{M})^{\perp}\).
The converse is the same implication read backwards, so
\eqref{v4-arith:w9-r4:eq:deficiency-range-defect-plus} follows.
The density criterion follows because a subspace has zero orthogonal
complement exactly when its closure is the full Hilbert space.
\end{proof}

\begin{corollary}[range density is graph density for \(M_t\)]
\label{v4-arith:w9-r4:cor:bulk-deficiency-graph-core}
\ClaimStatusProvedHere. For a subspace
\(\mathcal{C}\subset\mathcal{D}(M_t)\), the following are equivalent:
\[
\overline{\mathcal{C}}^{\,\|\cdot\|_{M_t}}=\mathcal{D}(M_t),
\qquad
\overline{(M_t-i)\mathcal{C}}^{\,L^2}=L^2(\mathbb R,dt),
\qquad
\overline{(M_t+i)\mathcal{C}}^{\,L^2}=L^2(\mathbb R,dt).
\]
Consequently, under the specified Mellin chart hypothesis,
\(\mathcal{C}_{\mathrm{res}}^{M}\) is a graph core for \(M_t\) if and
only if
\[
\mathcal{N}_{+}=0=\mathcal{N}_{-}.
\]
\end{corollary}

\begin{proof}
For \(\operatorname{Im}z\neq0\), the map
\[
M_{t}-z\colon \mathcal{D}(M_{t})_{\mathrm{graph}}\to L^{2}(\mathbb{R},dt)
\]
is boundedly invertible: its inverse is multiplication by
\((t-z)^{-1}\), which maps \(L^{2}\) into \(\mathcal{D}(M_{t})\) with
bounded graph norm. Hence a subspace is dense in the graph norm if
and only if its image under \(M_t-z\) is dense in \(L^2\). Apply this
with \(z=i\) and \(z=-i\), then use
\Cref{v4-arith:w9-r4:thm:deficiency-range-defect}.
\end{proof}

\begin{lemma}[rank-one graph-boundary model]
\label{v4-arith:w9-r4:lem:rank-one-graph-boundary}
\ClaimStatusProvedHere. Let
\(\psi(t)=(1+|t|)^{-1}\). Then
\(\psi\in L^{2}(\mathbb{R},dt)\) but
\(\psi\notin\mathcal{D}(M_t)\). Define
\[
\mathcal{C}_{\psi}
:=
\{\,f\in\mathcal{D}(M_t)\mid
\langle (M_t-i)f,\psi\rangle=0\,\}.
\]
Then \(\mathcal{C}_{\psi}\) is \(L^{2}\)-dense and graph-closed of
codimension one, but it is not graph-dense. Its range defects are
\[
\bigl((M_t-i)\mathcal{C}_{\psi}\bigr)^{\perp}
=\mathbb C\psi,
\qquad
\bigl((M_t+i)\mathcal{C}_{\psi}\bigr)^{\perp}
=
\mathbb C\,\frac{t+i}{t-i}\psi .
\]
\end{lemma}

\begin{proof}
The function \(\psi\) is square-integrable, while
\(t\psi(t)\) is not square-integrable. Since
\[
M_t-i\colon\mathcal{D}(M_t)_{\mathrm{graph}}\to L^2(\mathbb R,dt)
\]
is a bounded isomorphism, \(\mathcal{C}_{\psi}\) is the inverse image
of the closed hyperplane \(\psi^\perp\); hence it is graph-closed of
codimension one and its first range defect is \(\mathbb C\psi\).
The defining functional is not continuous for the \(L^2\)-topology on
\(\mathcal{D}(M_t)\), because such continuity would represent it by
an \(L^2\)-vector, equivalently would put \(\psi\) in
\(\mathcal{D}(M_t)\). Therefore its kernel is \(L^2\)-dense.

For the second range, write
\[
(M_t+i)(M_t-i)^{-1}
\]
as multiplication by \((t+i)/(t-i)\), a unitary multiplier on \(L^2\).
The orthogonal complement of its image of \(\psi^\perp\) is the span
of \((t+i)(t-i)^{-1}\psi\).
\end{proof}

\begin{lemma}[point conditions are graph-invisible]
\label{v4-arith:w9-r4:lem:point-conditions-graph-invisible}
\ClaimStatusProvedHere. Let \(\Gamma\subset\mathbb R\) be locally
finite. Then \(C_c^\infty(\mathbb R\setminus\Gamma)\) is graph-dense
in \(\mathcal{D}(M_t)\). In particular, a boundary condition that is
only pointwise at the zero ordinates cannot obstruct graph density for
the ordinary \(M_t\)-graph norm.
\end{lemma}

\begin{proof}
The graph norm of \(M_t\) is the \(L^2\)-norm for the measure
\((1+t^2)\,dt\). Locally finite sets have zero measure for this
weight. Given \(f\in\mathcal{D}(M_t)\), truncate \(f\) to a compact
interval, remove intervals of arbitrarily small weighted measure
around the finitely many points of \(\Gamma\) in that interval, and
then mollify. The approximants lie in
\(C_c^\infty(\mathbb R\setminus\Gamma)\) and converge in the graph
norm.
\end{proof}

\begin{definition}[Hadamard-supported graph functional]
\label{v4-arith:w9-r4:def:hadamard-supported-graph-functional}
\ClaimStatusDefinitional. Let
\[
\Gamma_{\xi}^{\mathrm{crit}}:=\{\,\gamma\in\mathbb R\mid
\tfrac12+i\gamma\text{ is a non-trivial zero of }\zeta\,\}
\]
with multiplicities forgotten. This is the real critical-line
ordinate support relevant to the Hilbert--P\'olya target; it is not an
assertion that the full Hadamard divisor has already been realised on
\(\mathbb R\). A continuous functional
\[
\beta\colon \mathcal{D}(M_t)_{\mathrm{graph}}\longrightarrow\mathbb C
\]
is \emph{Hadamard-supported on the zero ordinates} if
\[
\beta(\phi)=0
\qquad
\text{for every }
\phi\in C_c^\infty(\mathbb R\setminus\Gamma_{\xi}^{\mathrm{crit}}).
\]
This is the graph-norm formulation of the idea that the boundary
condition sees only point or finite-jet Hadamard residue data at the
zero ordinates.
\end{definition}

\begin{theorem}[no Hadamard residue functional in the ordinary graph dual]
\label{v4-arith:w9-r4:thm:no-hadamard-graph-dual}
\ClaimStatusProvedHere. The ordinary \(M_t\)-graph dual contains no
non-zero Hadamard-supported functional. Equivalently,
\[
\beta\in(\mathcal{D}(M_t)_{\mathrm{graph}})^{*},
\qquad
\operatorname{supp}\beta\subset\Gamma_{\xi}^{\mathrm{crit}}
\quad\Longrightarrow\quad
\beta=0
\]
in the sense of
\Cref{v4-arith:w9-r4:def:hadamard-supported-graph-functional}.
Thus there is no non-zero
\[
\beta_{\mathrm{Had}}\colon \mathcal{D}(M_t)_{\mathrm{graph}}\to\mathbb C
\]
whose graph-closed kernel is exactly a Hadamard-residue boundary
condition supported only at the zero ordinates. The statement also
rules out finite-jet residue data at those ordinates, if such data are
asked to be continuous for the ordinary \(M_t\)-graph topology.
\end{theorem}

\begin{proof}
The non-trivial zeros of \(\zeta\) form a discrete subset of the
critical strip, and only finitely many have ordinate in a bounded
interval. Hence \(\Gamma_{\xi}^{\mathrm{crit}}\) is locally finite as a subset of
\(\mathbb R\). By
\Cref{v4-arith:w9-r4:lem:point-conditions-graph-invisible},
\[
C_c^\infty(\mathbb R\setminus\Gamma_{\xi}^{\mathrm{crit}})
\]
is dense in \(\mathcal{D}(M_t)\) for the graph norm. If
\(\beta\) is continuous for the graph norm and vanishes on this dense
subspace, it vanishes on all of \(\mathcal{D}(M_t)\).

Equivalently, the graph norm is the \(L^2\)-norm for the weighted
measure \((1+t^2)\,dt\), so the graph dual is represented by weighted
\(L^2\)-functions, or, after the isometric resolvent chart
\[
M_t-i\colon \mathcal{D}(M_t)_{\mathrm{graph}}\xrightarrow{\sim}
L^2(\mathbb R,dt),
\]
by \(L^2\)-vectors. No non-zero \(L^2\)-vector is supported on the
locally finite set \(\Gamma_{\xi}^{\mathrm{crit}}\), which has Lebesgue measure zero.
\end{proof}

\begin{corollary}[Hadamard residue data is graph-invisible]
\label{v4-arith:w9-r4:cor:beta-had-no-go}
\ClaimStatusProvedHere. In the ordinary Mellin graph topology,
\(\beta_{\mathrm{Had}}\) does not exist as a non-zero graph-dual
boundary functional. Therefore a boundary condition whose only data are
point or finite-jet Hadamard residues at the locally finite real
ordinate set \(\Gamma_{\xi}^{\mathrm{crit}}\) is graph-invisible: it
cannot by itself obstruct graph density of
\(\mathcal{C}_{\mathrm{res}}^{M}\) in \(\mathcal{D}(M_t)\).
\end{corollary}

\begin{remark}[surviving branches after the graph-dual no-go]
\label{v4-arith:w9-r4:rem:beta-had-surviving-branches}
The preceding corollary is not an exhaustive boundary-state
classification. It eliminates only the ordinary point-supported graph-dual path. Assuming
the closed Tate-polar quotient and specified Mellin chart have been
constructed, a surviving operator obstruction must take a different form.
The known nonexclusive presentations are:
\begin{enumerate}
\item the exact Hadamard residue or finite-jet trace lives in a stronger
rigged, Paley--Wiener, de Branges, or reproducing-kernel topology
\(\Phi_{\xi}\subset\mathcal{D}(M_t)_{\mathrm{graph}}\), where point
evaluation is continuous;
\item the ordinary graph obstruction is nonlocal, represented under the
resolvent chart by an \(L^2\)-vector or closed \(L^2\)-subspace as in
\Cref{v4-arith:w9-r4:lem:rank-one-graph-boundary}, and is therefore
not point-residue data;
\item a physical Ward/BV--BFV argument supplies extra boundary phase
space not yet identified with
\(\mathcal{Q}_{\mathrm{res}}\); this is a physics residual, not an
operator-theoretic graph-dual boundary functional.
\end{enumerate}
These alternatives use different topologies and boundary spaces.
Identifying any one with a specified arithmetic Tate-positivity or
Burnol--de Branges construction requires a pairing isometry, an exact
operator-domain comparison, and the arithmetic trace identity.
No such identification follows from the list of alternatives.
The stronger-topology construction requires, at minimum, a continuous
Hadamard trace-and-jet map
\(\Phi_{\xi}\to\ell^{2}(\Gamma_{\xi}^{\mathrm{crit}},w)\), compatibility
with the post-Tate quotient and adjoint domains, and a proof that the
resulting closed condition is not a tautological insertion of the zero
divisor.
\end{remark}

\begin{remark}[answer of the calculation]
\label{v4-arith:w9-r4:rem:deficiency-calculation-answer}
The maximal multiplication operator on ordinary $L^2$ is
self-adjoint. For a restriction transported by the specified quotient
chart, essential self-adjointness asks exactly for graph density:
\[
\overline{\mathcal{C}_{\mathrm{res}}^{M}}^{\,\|\cdot\|_{M_{t}}}
\stackrel{?}{=}
\mathcal{D}(M_{t}),
\qquad
\|f\|_{M_{t}}^{2}:=\|f\|_{2}^{2}+\|t f\|_{2}^{2},
\]
in the quotient graph norm. A genuine boundary state must therefore
come from a graph-dual functional of the kind modeled in
\Cref{v4-arith:w9-r4:lem:rank-one-graph-boundary}, not from mere
point conditions at the zero ordinates, which are graph-invisible by
\Cref{v4-arith:w9-r4:lem:point-conditions-graph-invisible}. The direct
Hadamard-residue candidate \(\beta_{\mathrm{Had}}\) is ruled out in
the ordinary graph dual by
\Cref{v4-arith:w9-r4:thm:no-hadamard-graph-dual}.
\end{remark}

\subsection{Graph Norm and Arithmetic Determinant}
\label{v4-arith:w9-r4:pdnf:execution-attempt}

The continuous Mellin operator has absolutely continuous spectrum
$\mathbb R$. A positive atomic model takes its zero divisor and
weights as input. A quotient constructed from arithmetic data would
require its own positive pairing, graph domain, and trace comparison.
The following calculations distinguish the first two constructions
and the conditions required to compare either with the third.

\begin{lemma}[continuous Mellin-line graph closure]
\label{v4-arith:w9-r4:lem:meyer-line-graph-closure}
\ClaimStatusProvedHere. Let
\[
M_{t}^{0}:C_{c}^{\infty}(\mathbb{R})\subset L^{2}(\mathbb{R},dt)
\longrightarrow L^{2}(\mathbb{R},dt),
\qquad
(M_{t}^{0}f)(t)=t f(t).
\]
Then \(M_{t}^{0}\) is essentially self-adjoint. Its graph-norm
closure is multiplication by \(t\) on
\[
\mathcal{D}(M_{t})
\;=\;
\{\,f\in L^{2}(\mathbb{R},dt)\mid t f(t)\in L^{2}(\mathbb{R},dt)\,\}.
\]
The spectrum is the absolutely continuous spectrum
\(\mathrm{Spec}(M_{t})=\mathbb{R}\).
\end{lemma}

\begin{proof}
The operator \(M_{t}\) with the displayed maximal domain is
self-adjoint by the spectral theorem for multiplication operators.
The core \(C_{c}^{\infty}(\mathbb{R})\) is dense in the graph norm:
truncate \(f\) by compactly supported cutoffs and then mollify; both
\(f\) and \(t f\) converge in \(L^{2}\). Hence the closure of
\(M_{t}^{0}\) is \(M_{t}\). Since \(M_{t}\) is multiplication by the
coordinate on \(L^{2}(\mathbb{R},dt)\), its spectrum is the essential
range of \(t\), namely \(\mathbb{R}\), with absolutely continuous
spectral measure.
\end{proof}

\begin{proposition}[why the classical closure is insufficient]
\label{v4-arith:w9-r4:prop:classical-closure-insufficient}
\ClaimStatusProvedHere. The closure of
\Cref{v4-arith:w9-r4:lem:meyer-line-graph-closure} does not solve
the spectral-flow problem. Its graph norm is complete and its
operator is self-adjoint, but its spectrum is continuous
\(\mathbb{R}\), not a divisor with zero multiplicities of
\(\xi_{\mathbb R}\).
Consequently it cannot satisfy
\[
\det\nolimits_{\infty}
\bigl(s\cdot\mathrm{id}-({\tfrac12}\mathrm{id}+iM_{t})\bigr)
=\xi_{\mathbb{R}}(s)
\]
as a zero-eigenvalue determinant.
\end{proposition}

\begin{proof}
The determinant identity in
\Cref{v4-arith:w9-r4:eq:closure-det} has divisor
\(\mathrm{div}\,\xi_{\mathbb R}\). The multiplication operator
\(M_{t}\) on \(L^{2}(\mathbb{R},dt)\) has Lebesgue spectral measure on
the line, not this determinant divisor. Thus the closure obtained by
ordinary Mellin-line graph completion has continuous spectral
type and no point eigenvalues. This statement concerns the continuous
Hilbert space. Meyer's global nuclear bornological representation
has the separate character and multiplicity theorem recorded in
Theorem~\ref{v4-arith:w9-r4:thm:meyer-intertwiner}; it is not
identified with this Hilbert closure.
\end{proof}

\begin{proposition}[formal zero-space closure is tautological]
\label{v4-arith:w9-r4:prop:formal-zero-closure-tautological}
\ClaimStatusProvedHere. On the formal zero Hilbert space
\(\mathcal{H}_{\xi}^{\mathrm{formal}}\) of
\Cref{v4-arith:w9-r4:def:target-data}, the diagonal ordinate
operator has graph-norm closure. Its projected operator
\(\Theta_{\xi}^{\mathrm{proj}}\) satisfies the regularised
determinant identity for \(\xi_{\mathbb R}\) only if the original
zero divisor already lies on the critical line. This proves no
non-formal spectral-flow closure.
\end{proposition}

\begin{proof}
The finite-support diagonal operator on
\(\bigoplus_{\rho}\mathbb{C}^{m_{\rho}}\) closes in the graph norm to
the maximal diagonal multiplication operator by the real numbers
\(t_{\rho}=\mathrm{Im}(\rho)\); its closure is self-adjoint. The
operator \(\Theta_{\xi}^{\mathrm{proj}}=\tfrac12+iH_{\mathrm{ord}}^{\mathrm{formal}}\)
has spectral divisor
\(\{\tfrac12+it_{\rho}\}\). Hadamard factorisation identifies its
regularised determinant with \(\xi_{\mathbb R}\) only when
\(\beta_{\rho}=1/2\) for every \(\rho\). Otherwise it is the
determinant of the projected divisor, not of \(\xi_{\mathbb R}\).
The construction has therefore used the zero divisor and then
discarded its real-part defect. It supplies no Bost--Connes, Meyer,
Connes, Deninger, or chiral-homological source for the Hilbert space
or the Weil--Connes trace; it is the formal target explicitly excluded
by \Cref{v4-arith:w9-r4:rem:why-formal}.
\end{proof}

\begin{theorem}[the Dirac criterion on the fixed quotient]
\label{v4-arith:w9-r4:thm:execution-result}
\ClaimStatusProvedHereConditional. Suppose the positive quotient
Hilbert space and densely defined symmetric operator
$T_0=H_0^\circ$ have been constructed as above.
Fix the determinant convention of
Theorem~\ref{v4-arith:w9-r4:thm:polyakov-dirac-normal-form}.
Then the following statements are equivalent:
\begin{enumerate}
\item $T_0$ is essentially self-adjoint and its closure satisfies
\eqref{v4-arith:w9-r4:eq:closure-det};
\item $\ker(T_0^*-i)=0=\ker(T_0^*+i)$, and the self-adjoint closure
has the same determinant identity;
\item the Dirac double of this same $T_0$ has the Dirac closure
of Definition~\ref{v4-arith:w9-r4:def:dirac-double}.
\end{enumerate}
Under a specified unitary graph-core comparison with an arithmetic
Hilbert--P\'olya operator, these conditions transfer to that operator.
They do not construct this comparison or an arithmetic trace map.
\end{theorem}
\begin{proof}
The equivalence of the first two statements is the deficiency-index
criterion for a densely defined symmetric operator.
The equivalence with the third is the fixed-domain Dirac
conjugacy proved in
Theorem~\ref{v4-arith:w9-r4:thm:polyakov-dirac-normal-form}.
The additional unitary transports closures, adjoints, and the
determinant under the hypotheses of
Proposition~\ref{v4-arith:mc:core-unitary} and
Corollary~\ref{v4-arith:mc:determinant}.
\end{proof}

\begin{remark}[the arithmetic significance of the divisor]
\label{v4-arith:w9-r4:rem:failed-execution-record}
A self-adjoint realization with the complete xi zero divisor
implies critical-line placement, because its shifted eigenvalues
have real part $1/2$. A graph closure on the ordinary Mellin line
has continuous spectral type. A Hilbert space built by prescribing
the zeros requires an additional arithmetic pairing and trace
comparison. The preceding fixed-core equivalence distinguishes
these constructions and preserves the original arithmetic question.
\end{remark}

\section{Bibliography}
\label{v4-arith:w9-r4:biblio}

The references below are grouped by the four mathematical
constructions used above.

\paragraph{Hilbert-P\'olya foundational literature.}
\begin{description}
\item[P\'olya.] G.~P\'olya, correspondence with Landau (1914),
transcribed in Pintz and S\'ark\"ozy, \emph{On the work of P\'olya in
analytic number theory}, Stud.~Sci.~Math.~Hungar.~26 (1991), 201--226.
\item[Hilbert.] D.~Hilbert, G\"ottingen seminar notes (1914, 1916),
recorded by E.~Landau, \emph{Handbuch der Lehre von der Verteilung der
Primzahlen}, Teubner (1909), 2nd ed.~1953.
\item[Berry--Keating.] M.~V.~Berry and J.~P.~Keating, \emph{The
Riemann zeros and eigenvalue asymptotics}, SIAM Review 41 (1999),
236--266.
\item[Montgomery.] H.~L.~Montgomery, \emph{The pair correlation of zeros
of the zeta function}, Proc.~Sympos.~Pure Math.~24 (1973), 181--193.
\end{description}

\paragraph{Adelic and idele-class operator constructions.}
\begin{description}
\item[Connes.] A.~Connes, \emph{Trace formula in noncommutative
geometry and the zeros of the Riemann zeta function}, Selecta
Math.~(N.~S.)~5 (1999), 29--106; arXiv:math/9811068. \S IV for the
absorption-trace conjecture; \S V for the Weil-positive trace
identity.
\item[Connes--Marcolli.] A.~Connes and M.~Marcolli, \emph{Noncommutative
Geometry, Quantum Fields, and Motives}, AMS Colloq.~Publ.~55 (2008).
Ch.~3 for KMS states; Ch.~4 for rigidity; Ch.~5 for the adele class
space.
\item[Meyer.] R.~Meyer, \emph{On a representation of the idele class
group related to primes and zeros of \(L\)-functions}, Duke Math.~J.
127 (2005), 519--595; arXiv:math/0311468v3. Theorem~5.8,
p.~43, gives the character formula. Lemma~5.10, p.~45, concerns
Fourier--Laplace transforms. Theorem~5.11, p.~47, gives spectral
multiplicities. Page locators refer to the arXiv text.
\item[Deninger.] C.~Deninger, \emph{Motivic \(L\)-functions and
regularized determinants}, Proc.~Sympos.~Pure Math.~55.1 (1994),
707--743; \emph{Some analogies between number theory and dynamical
systems on foliated spaces}, Proc.~ICM~Berlin~1 (1998), 163--186.
\end{description}

\paragraph{Mellin and Hadamard classical analysis.}
\begin{description}
\item[Mellin.] H.~Mellin, \emph{Die Dirichletschen Reihen, die
zahlentheoretischen Funktionen und die unendlichen Produkte von
endlichem Geschlecht}, Acta Soc.~Sci.~Fenn.~31, no.~2 (1904).
\item[Hadamard.] J.~Hadamard, \emph{\'Etude sur les propri\'et\'es
des fonctions enti\`eres et en particulier d'une fonction consid\'er\'ee
par Riemann}, J.~Math.~Pures Appl.~9 (1893), 171--215.
\item[Titchmarsh.] E.~C.~Titchmarsh, \emph{Introduction to the Theory
of Fourier Integrals}, 2nd ed., Oxford (1948); \emph{The Theory of
the Riemann Zeta-Function}, 2nd ed., revised D.~R.~Heath-Brown, Oxford
(1986).
\item[von Mangoldt.] H.~von Mangoldt, \emph{Zu Riemann's Abhandlung
``\"Uber die Anzahl der Primzahlen unter einer gegebenen Gr\"osse''},
J.~Reine Angew.~Math.~114 (1895), 255--305.
\item[Edwards.] H.~M.~Edwards, \emph{Riemann's Zeta Function},
Academic Press (1974). \S2.7 for Hadamard product; \S3 for the
explicit formula.
\item[Patterson.] S.~J.~Patterson, \emph{An Introduction to the Theory
of the Riemann Zeta-Function}, Cambridge (1988). Ch.~3 for the Weil
explicit formula.
\end{description}

\paragraph{Bost--Connes and GNS modular theory.}
\begin{description}
\item[Bost--Connes.] J.-B.~Bost and A.~Connes, \emph{Hecke algebras,
type~III factors and phase transitions with spontaneous symmetry
breaking in number theory}, Selecta Math.~1 (1995), 411--457.
\item[Takesaki.] M.~Takesaki, \emph{Theory of Operator Algebras~II},
Springer (2003). Ch.~VIII for modular automorphism groups.
\item[Julia.] B.~Julia, \emph{Statistical theory of numbers}, Springer
Proc.~Phys.~47 (1990), 276--293. Primon-gas interpretation of the
Bost--Connes spectrum.
\end{description}

\paragraph{First-order and Dirac normal-form sources.}
\begin{description}
\item[Dirac.] P.~A.~M.~Dirac, \emph{The quantum theory of the
electron}, Proc.~Roy.~Soc.~Lond.~A 117 (1928), 610--624, and
118 (1928), 351--361. Source for the first-order square-root
principle used here only as a normal-form discipline.
\item[Reed--Simon.] M.~Reed and B.~Simon, \emph{Methods of Modern
Mathematical Physics I: Functional Analysis}, Academic Press (1980).
Ch.~VIII for adjoints, closures, block operators, and the
self-adjointness criterion used in
\Cref{v4-arith:w9-r4:lem:dirac-double-detects-sa}.
\end{description}

\paragraph{Vol~IV earlier sources used here.}
\begin{description}
\item \Cref{v4-arith:w5-b:hp-operator}: Bost--Connes GNS and
log-modular operator construction (Ch.~\ref{v4-arith:w5-b:hp-operator}).
\item \Cref{v4-arith:rh:hyp:HPAr}: Hilbert--P\'olya hypothesis
(Ch.~\ref{v4-arith:rh-reformulation}).
\item \Cref{v4-arith:HPAr-vbconverse}: H-P\(^{\mathrm{Ar}}\) equivalence bundle, approach~2
(Ch.~\ref{v4-arith:HPAr-vbconverse}).
\item \Cref{v4-arith:hpar-det}: determinant realisation of H-P\(^{\mathrm{Ar}}\), approach~3
(Ch.~\ref{v4-arith:hpar-det}).
\item \Cref{v4-arith:w7-j:ATP-adelic-scattering}: adelic-scattering
formulation (Ch.~\ref{v4-arith:w7-j:ATP-adelic-scattering}).
\item \Cref{v4-arith:w8-a:deninger-chiral-homology}: Deninger-chiral
triple identification (Ch.~\ref{v4-arith:w8-a:deninger-chiral-homology}).
\item \Cref{v4-arith:3d-acs-E1bar}: 3D arithmetic Chern-Simons
(Ch.~\ref{v4-arith:3d-acs-E1bar}).
\end{description}
